\section{Shared results III: hermitian curvature, growth and divisibility bounds}\label{sec:pf:sh-curv}

\providecommand{\shcurvDs}{d_\sigma}
\providecommand{\shcurvDt}{d_\tau}
\providecommand{\shcurvGap}{h_0}
\providecommand{\shcurvRam}{\mathrm{ram}}
\providecommand{\shcurvPP}{\mathsf P}
\providecommand{\shcurvNm}{\mathfrak n^-}

This section proves fourteen results of the formal development \cite{LeanSource2026} which are used by several of the
preceding sections of Part~II. They fall into four independent groups. Some are classical theorems or formalized forms of
them (the extremal projector of Asherova, Smirnov and Tolstoy, Harish-Chandra's formula for the formal degree of discrete
series); for these we state precisely what the formal statement asserts and how it relates to the
published one. Others are re-packagings which get a short proof. The remaining ones rest on long formal proofs; for them we
give the statement, the complete deduction from the main ingredients, and the statements of the ingredients, some of which
are only cited with their content.

\begin{itemize}
\item \Cref{sec:pf:sh-curv:vol}: volumes, formal degrees and slopes (used in
\cref{sec:pf:stageU,sec:pf:top-shape,sec:pf:open0}): volume
ratios of hermitian lattices are mass ratios (\cref{pf:sh-curv:volmass}), two-sided bounds for the formal degrees of discrete
series of $\Uni(p,q)$ in large gap (\cref{pf:sh-curv:dsig}), and the slope line of the ramified lattice
(\cref{pf:sh-curv:slopeline}).
\item \Cref{sec:pf:sh-curv:ep}: the extremal projector and the growth bounds (used in \cref{sec:pf:linkG}): the extremal
projector statement (\cref{pf:sh-curv:E}), the tool statements (\cref{pf:sh-curv:tools}) and the threshold functions of the
first theorem (\cref{pf:sh-curv:firstthm}).
\item \Cref{sec:pf:sh-curv:div}: divisibility (used in \cref{sec:pf:reckernel,sec:pf:top-shape,sec:pf:onto3}): the orbit
letter bounds with constants $(11,1)$ (\cref{pf:sh-curv:split}) and exponential gauges (\cref{pf:sh-curv:expgauge}).
\item \Cref{sec:pf:sh-curv:arch}: archimedean inputs at built places (used in
\cref{sec:pf:thmS-carrier,sec:pf:thmS-leaves,sec:pf:thmS}): convergence of orbital integrals (\cref{pf:sh-curv:orb}), top
lifts (\cref{pf:sh-curv:toplift}), chart equations (\cref{pf:sh-curv:charteq}), the gauge transfer, whose central step is a
hermitian curvature computation (\cref{pf:sh-curv:gentransfer}), the head-top identity (\cref{pf:sh-curv:headtop}), and
Theorem P for compactly supported test functions (\cref{pf:sh-curv:actualcpt}), the largest result of the section.
\end{itemize}

% ===========================================================================
\subsection{Volumes, formal degrees and slopes}\label{sec:pf:sh-curv:vol}

The three results of this subsection are used in the proofs of the family properties (G1)--(G11) of
\cref{sec:pf:top-shape} (and of their second-layer form in \cref{sec:pf:open0}) and of the genus-mass rows of
\cref{sec:pf:stageU}. Throughout, $D$ is a datum with host CM
field $E$, $F\subset E$ its maximal totally real subfield, $f=[F:\Q]$, and
$v=|\mathrm{disc}\,E|^{1/[E:\Q]}/2$ is the invariant of \cref{prop:families:mass}. For $M\ge1$ we write
\[
L_\eta(M)=\prod_{\mathfrak P}e_M(\mathfrak P)
\]
for the product of \cref{prop:families:mass}, over the prime ideals $\mathfrak P$ of $\OO_E$; for data it is a
positive real number (\lref{Formalism.MainNotions.StandardLattice.lVal_pos}, which states $L_\eta(M)>0$ for $M\ge2$).

\subsubsection{Volume ratios and mass ratios}

\begin{definition}[Volume ratios and mass ratios]\label{pf:sh-curv:vrmr}
A \emph{hermitian space} $W$ over $E$ of rank $n$ consists of a non-degenerate hermitian form $\langle\ ,\ \rangle$ on an
$n$-dimensional $E$-vector space (with a frame and a unimodular base lattice) and a signature $(P_\tau,Q_\tau)$ for every
embedding $\tau\colon E\to\C$; it is \emph{well signed} if $(P_\tau,Q_\tau)$ is the signature of the Gram matrix of the
frame under $\tau$, for every $\tau$. The datum attaches a volume $\vol(W)>0$ to every hermitian space. Consider the
following two statements about $D$, for well-signed $W$ of rank $n\ge3$ with signature $(P,Q)$ at every $\tau$ and a sign
$\pm$ with $P\ge1$ (for $+$), resp.\ $Q\ge1$ (for $-$):
\begin{itemize}
\item[(VR$_D$)] there are $x\in W$ with $\langle x,x\rangle=\pm1$, a well-signed hermitian space $W'$ of signature
$(P-1,Q)$ (resp.\ $(P,Q-1)$) at every $\tau$ and an isometric embedding $\iota\colon W'\to W$ with image $x^\perp$, such
that
\begin{equation}\label{eq:pf:sh-curv:VR}
\frac{\vol(W')}{\vol(W)}=\Bigl(\frac{(2\pi)^n}{(n-1)!}\Bigr)^f\bigl((2v)^f\bigr)^{\frac12-n}L_\eta(n)^{-1};
\end{equation}
\item[(MR$_D$)] the same, with \eqref{eq:pf:sh-curv:VR} replaced by
\begin{equation}\label{eq:pf:sh-curv:MR}
\frac{\vol(W)}{\vol(W')}=\frac{\bigl((n-1)!\,(v/\pi)^n\bigr)^f}{\sqrt{(2v)^f}}\,L_\eta(n),
\end{equation}
the right-hand side of the mass ratio of \cref{prop:families:mass} (with $d=f$ and $M=n$).
\end{itemize}
\formal{Formalism.MainNotions.StandardLattice.VolumeRatio₂}
\formal{Formalism.MainNotions.StandardLattice.Conclusion2}
\end{definition}

(VR$_D$) is the volume ratio (VR) of \cref{sec:pf:stageU}, specialised to the datum $D$.

\begin{lemma}[The global arithmetic identity]\label{pf:sh-curv:gai}
Let $f,M\ge1$ be integers and $v,L>0$ real numbers, and put $d_F=(2v)^f$ and $c_M=(2\pi)^M/(M-1)!$. Then
\[
\Bigl(c_M^{\,f}\,d_F^{\,\frac12-M}\,L^{-1}\Bigr)^{-1}=\frac{\bigl((M-1)!\,(v/\pi)^M\bigr)^f}{\sqrt{d_F}}\,L .
\]
\formallabel{Formalism.HermitianLattices.GenusMassVolumes.global_arithmetic_identity}
\end{lemma}

\begin{proof}
Put $A=c_M$ and $B=(M-1)!\,(v/\pi)^M$; both are positive, and $AB=(2\pi)^M(v/\pi)^M=(2v)^M$, so
$A^fB^f=(2v)^{fM}=d_F^{\,M}$. Since $d_F^{\,\frac12-M}=\sqrt{d_F}/d_F^{\,M}$, the left-hand side equals
$L\,d_F^{\,M}/(A^f\sqrt{d_F})=L\,B^f/\sqrt{d_F}$.
\end{proof}

\begin{lemma}[Volume ratios are mass ratios]\label{pf:sh-curv:volmass}
For every datum $D$, \textup{(VR$_D$)} implies \textup{(MR$_D$)}.
\formallabel{Formalism.MainNotions.StandardLattice.volumeMass_of_unfolding}
\end{lemma}

\begin{proof}
Let $W,P,Q,\pm$ be as in (MR$_D$), and take $x,W',\iota$ as provided by (VR$_D$); all conditions of (MR$_D$) except
\eqref{eq:pf:sh-curv:MR} are the same as in (VR$_D$). Both volumes are positive, so
$\vol(W)/\vol(W')$ is the inverse of the right-hand side of \eqref{eq:pf:sh-curv:VR}. Since $n\ge3$, $v>0$
(\lref{Formalism.MainNotions.StandardLattice.host_v_pos}, which states that the root discriminant of $E$ is positive)
and $L_\eta(n)>0$, \cref{pf:sh-curv:gai} with $M=n$ and $L=L_\eta(n)$ identifies this inverse with the right-hand side of
\eqref{eq:pf:sh-curv:MR}. Finally, the invariants of $E$ used in \cref{prop:families:mass} are those of the datum: the
number of infinite places of $E$ is $f$ and $|\mathrm{disc}\,E|^{1/(2f)}/2=v$
(\lref{Formalism.MainNotions.StandardLattice.mass_ratio_eq}, which states that the right-hand side of the mass ratio
computed from $E$ equals the expression \eqref{eq:pf:sh-curv:MR} computed from the datum).
\end{proof}

\subsubsection{Formal degrees of discrete series in large gap}

\begin{definition}[Discrete series of $\Uni(p,q)$]\label{pf:sh-curv:dsdef}
Let $N\ge2$ and let $\eps\in\{\pm\}^N$ be a sign vector with $p\ge1$ plus signs and $q\ge1$ minus signs. Let
$U_\eps\cong\Uni(p,q)$ be the unitary group of the hermitian form $\mathrm{diag}(\eps)$ on $\C^N$, $K_\eps\cong\Uni(p)\times
\Uni(q)$ the subgroup of its elements commuting with $\mathrm{diag}(\eps)$, and $\mu_\eps$ the reference Haar measure on
$U_\eps$ of the development. Put
\begin{gather*}
\rho_c(a)=\tfrac12\bigl(\#\{b>a:\eps_b=\eps_a\}-\#\{b<a:\eps_b=\eps_a\}\bigr),\\
\rho_n(a)=\tfrac12\bigl(\#\{b>a:\eps_b\ne\eps_a\}-\#\{b<a:\eps_b\ne\eps_a\}\bigr).
\end{gather*}
A \emph{Harish-Chandra parameter} is a strictly decreasing $\lambda\in\R^N$ with $\lambda_a-\frac{N-1}2\in\Z$ for all
$a$; it has \emph{gap at least $h$} if $|\lambda_a-\lambda_b|\ge h$ whenever $\eps_a\ne\eps_b$. We write
$|\lambda|=\max_{a,b}|\lambda_a-\lambda_b|$ and
\[
\dim V(\lambda)=\prod_{a<b}\frac{\lambda_a-\lambda_b}{b-a},\qquad
\shcurvDt(\eps,\lambda)=\frac{\prod_{a<b,\ \eps_a=\eps_b}\bigl((\lambda+\rho_n)_a-(\lambda+\rho_n)_b\bigr)}
{\prod_{j<p}j!\,\prod_{j<q}j!}.
\]
A \emph{discrete series representation with parameter $\lambda$} is an irreducible unitary representation $\sigma$ of
$U_\eps$ on a non-zero Hilbert space $H$ whose matrix coefficients are square integrable, which contains the $K_\eps$-type
of highest weight $\Lambda=\lambda+\rho_n-\rho_c$, and such that $\Lambda$ minimises $\|\mu+2\rho_c\|$ among the highest
weights $\mu$ of the $K_\eps$-types of $\sigma$. A number $d>0$ is a \emph{formal degree} of $\sigma$ with respect to a Haar
measure $\mu$ if $\int\langle\sigma(g)v,w\rangle\overline{\langle\sigma(g)v',w'\rangle}\,d\mu(g)=d^{-1}\langle v,v'\rangle
\overline{\langle w,w'\rangle}$ for all $v,w,v',w'\in H$ (Schur orthogonality).
\formal{Formalism.QuantumDeterminants.IsHCParam}
\formal{Formalism.QuantumDeterminants.GapAtLeast}
\formal{Formalism.MainNotions.StandardLattice.dimV}
\formal{Formalism.MainNotions.StandardLattice.dTau}
\formal{Formalism.MainNotions.S.IsDS}
\formal{Formalism.MainNotions.S.IsFdeg}
\end{definition}

Thus $\dim V(\lambda)$ is Weyl's dimension of the irreducible representation of $\GL_N$ with infinitesimal character
$\lambda$, $\shcurvDt(\eps,\lambda)$ is Weyl's dimension of the irreducible representation of $K_\eps$ of highest weight
$\Lambda$ (note $\Lambda+\rho_c=\lambda+\rho_n$), and $\Lambda$ is Vogan's minimal $K$-type.

\begin{definition}[Degree data of a datum]\label{pf:sh-curv:dsties}
A datum $D$ provides a signature vector $\eps=\eps(D)$ of length $N=r$ with $(p,q)=(p_Y,q_Y)$ the signature of its base
space, for every \emph{row type} $\rho$ a Harish-Chandra parameter $\lambda(\rho)\in\R^N$, a scale $s>0$ for the Haar
measure of $U_\eps$, and numbers $\shcurvDs(\rho)$. The \emph{ties} of a member $D$ of a family include the
\emph{Haar scale tie} ($s\,\mu_\eps$ is the Haar measure of the development's convention for $\Uni(p,q)$) and the
\emph{formal degree tie} (if $\sigma$ is a discrete series representation with parameter $\lambda(\rho)$ and $d$ is a formal
degree of $\sigma$ with respect to $s\mu_\eps$, then $d=\shcurvDs(\rho)$). For $0<\eps_0$ put
\[
\shcurvGap(\eps_0)=8192\,\min(p,q)\,(pq)^3\,\frac{\log N}{\eps_0^2}.
\]
\formal{Formalism.BasicNotions.cdScaleTie}
\formal{Formalism.BasicNotions.dsigTie}
\formal{Formalism.MainNotions.StandardLattice.spectralGapLowerBound}
\end{definition}

\begin{theorem}[Two-sided bounds for $\shcurvDs$]\label{pf:sh-curv:dsig}
Let $\mathcal F$ be a family. Then eventually every member $D$ of $\mathcal F$ satisfies: for all $\eps_0$ with
$e^{-N}\le\eps_0<\frac12$ and every row type $\rho$ whose parameter $\lambda=\lambda(\rho)$ has gap at least
$\shcurvGap(\eps_0)$,
\[
(1-\eps_0)\dim V(\lambda)\le\shcurvDs(\rho)\le4\,\shcurvDt(\eps,\lambda)\,(3pq\,|\lambda|)^{pq}.
\]
In fact this holds for every datum $D$ satisfying the Haar scale tie and the formal degree tie, hence for all members of
every family.
\formal{Formalism.MainNotions.StandardLattice.SigmaDegreeTwoSidedBoundOfGap}
\formallabel{Formalism.MainNotions.GappedMultiplicityBounds.explicit_bound_family}
\formallabel{Formalism.MainNotions.GappedMultiplicityBounds.theorem_three_family_of_gapped_multiplicity}
\formallabel{Formalism.MainNotions.GappedMultiplicityBounds.theorem_three_rows_of_real_place}
\formallabel{Formalism.MainNotions.GappedMultiplicityBounds.theorem_three_of_weyl_gapped_multiplicity}
\formallabel{Formalism.MainNotions.GappedMultiplicityBounds.theorem_three_of_gap_bound}
\formallabel{Formalism.MainNotions.GappedMultiplicityBounds.finrank_lowK_eqG}
\formallabel{PI.UnitaryGroupMassFormula.HCFdegRGap.gap_big}
\formallabel{PI.UnitaryGroupMassFormula.HCFdegRGap.dsigEx_at}
\end{theorem}

The proof uses three results on discrete series of $\Uni(p,q)$ which are proved elsewhere in the development:
\begin{itemize}
\item (Existence) If $\lambda$ is a Harish-Chandra parameter with gap $>\max(2N,8N\min(p,q))$, there is a discrete series
representation with parameter $\lambda$ (\lref{ReciprocityTheory.ChartRestriction.DS.G2.discrete_series_exists_large_gap});
every discrete series representation has a formal degree with respect to $\mu_\eps$
(\lref{Formalism.MainNotions.BridgeCSchur.fdeg_of_DS}). These are classical
\cite{HarishChandra1966}.
\item (Minimal $K$-type) If $\lambda$ has gap $>Npq+1$, the $K_\eps$-type of highest weight $\Lambda$ occurs exactly once in
a discrete series representation with parameter $\lambda$
(\lref{Formalism.MainNotions.BesselDiscreteSeries.lowest_K_type_mult_one_of_gaps}); this is a special case of
Blattner's formula, proved by Hecht and Schmid \cite{HechtSchmid1975}.
\item (Weyl) If $S$ is an irreducible $K_\eps$-type of highest weight $\mu$ in a unitary representation of $K_\eps$, then
$\dim S=\prod_{a<b,\ \eps_a=\eps_b}\bigl((\mu+\rho_c)_a-(\mu+\rho_c)_b\bigr)\big/\bigl(\prod_{j<p}j!\prod_{j<q}j!\bigr)$
(\lref{weylForK_closed}); this is Weyl's dimension formula for $\Uni(p)\times\Uni(q)$.
\end{itemize}

\begin{lemma}[Harish-Chandra's formal degree in large gap]\label{pf:sh-curv:hc}
Let $\lambda$ be a Harish-Chandra parameter with gap $h>32Npq+2$, let $\sigma$ be a discrete series representation with
parameter $\lambda$, and let $s\mu_\eps$ be a Haar measure satisfying the Haar scale tie. Then for every $w\ne0$,
\[
\frac{\|w\|^4}{\int_{U_\eps}|\langle\sigma(g)w,w\rangle|^2\,d(s\mu_\eps)(g)}=\prod_{a<b}\frac{|\lambda_a-\lambda_b|}{b-a}.
\]
\formallabel{PI.UnitaryGroupMassFormula.HCFdegRGap.hCFdegRGap_holds}
\formallabel{PI.UnitaryGroupMassFormula.HCFdegRGap.fdeg_cd_gap}
\formallabel{PI.UnitaryGroupMassFormula.HCFdegRGap.thrStar_lt}
\formallabel{PI.UnitaryGroupMassFormula.HCFdegRGap.muCD_eq_smul}
\formallabel{Formalism.MainNotions.StandardLattice.formal_degree_of_eq}
\end{lemma}

\begin{proof}
By the Haar scale tie, $s\mu_\eps$ is the measure $\mu^{\mathrm{CD}}_\eps$ for which the development proves
Harish-Chandra's formula: both are the push-forward of the reference measure of $\Uni(p,q)$ under a reordering of the
coordinates putting the signs of $\eps$ in standard position (for $\mu^{\mathrm{CD}}_\eps$ this is
\lref{ReciprocityTheory.ChartRestriction.DS.muCD_eq_cd_vol}, for $s\mu_\eps$ it is the tie). The gap hypothesis implies
$h>\max\bigl(Npq+1,\,16N\min(p,q)+8\sqrt{2\min(p,q)}\bigr)+1$: indeed $\min(p,q)\le pq$, $pq\le Npq$ and
$\sqrt{2m}\le2m$ for $m\ge1$, so both entries of the maximum are at most $32Npq+1$. Under this hypothesis, the formal
degree of $\sigma$ with respect to $\mu^{\mathrm{CD}}_\eps$ is $\prod_{a<b}|\lambda_a-\lambda_b|/(b-a)$
(\lref{ReciprocityTheory.ChartRestriction.DS.fdCons_cd}, proved from the existence of an $L^2$ eigenfunction of the
Casimir operator transforming under the minimal $K_\eps$-type,
\lref{ReciprocityTheory.ChartRestriction.DS.Lap.Base.psiFor}); this is Harish-Chandra's formula for the formal
degree \cite{HarishChandra1966} with the normalisation of the measure built into $\mu^{\mathrm{CD}}_\eps$. Finally, if
$d$ is a formal degree, Schur orthogonality with $v=w=v'=w'$ gives $\int|\langle\sigma(g)w,w\rangle|^2=\|w\|^4/d$, so the
left-hand side equals $d$.
\end{proof}

\begin{lemma}\label{pf:sh-curv:fdegle}
For every sign vector $\eps$ with $p$ plus and $q$ minus signs and every Harish-Chandra parameter $\lambda$,
$\shcurvDt(\eps,\lambda)>0$ and
\[
\prod_{a<b}\frac{|\lambda_a-\lambda_b|}{b-a}\le\shcurvDt(\eps,\lambda)\,|\lambda|^{pq}.
\]
\formallabel{Formalism.L2.GlobalTheoryTheory.GenusMassVolumes.fdegHC_le_dTau}
\formallabel{Formalism.MainNotions.StandardLattice.formalDegreeFactor_pos_of_strictAnti}
\end{lemma}

\begin{proof}
Since $\lambda$ is strictly decreasing, $\lambda_a-\lambda_b>0$ for $a<b$. For $a<b$ with $\eps_a=\eps_b$ a direct count
gives
\[
(\lambda+\rho_n)_a-(\lambda+\rho_n)_b=\lambda_a-\lambda_b+\#\{c:a<c<b,\ \eps_c\ne\eps_a\}\ge\lambda_a-\lambda_b>0,
\]
so $\shcurvDt(\eps,\lambda)>0$; and for $a<b$ with $\eps_a\ne\eps_b$ trivially $\lambda_a-\lambda_b\le|\lambda|$. Splitting the product over $a<b$ according
to $\eps_a=\eps_b$ or not, the numerator $\prod_{a<b}(\lambda_a-\lambda_b)$ is at most
$\prod_{a<b,\eps_a=\eps_b}((\lambda+\rho_n)_a-(\lambda+\rho_n)_b)\cdot|\lambda|^{m}$, where $m$ is the number of pairs $a<b$
with $\eps_a\ne\eps_b$; $m\le pq$. If $N\ge2$ then $|\lambda|\ge1$ (the differences $\lambda_a-\lambda_b$ are non-zero
integers), so $|\lambda|^m\le|\lambda|^{pq}$; if $N\le1$ then $m=pq=0$. For the denominators,
$\prod_{a<b}(b-a)=\prod_{b<N}b!=\prod_{j<N}j!$, and $\prod_{j<p}j!\prod_{j<q}j!\le\prod_{j<p+q}j!$ since
$j!\le(p+j)!$. As the remaining factors are non-negative, the claim follows.
\end{proof}

\begin{proof}[Proof of \cref{pf:sh-curv:dsig}]
Formally the statement is that there is $x_0$ such that the bound holds at all members $D(x,i)$ with $x\ge x_0$; we take
$x_0=0$ and prove the bound for every datum $D$ satisfying the two ties, which the members of a family do (they are part
of the ties of the family). Let $\eps_0,\rho$ be as in the statement, $\lambda=\lambda(\rho)$ and $h_0=\shcurvGap(\eps_0)$.

\emph{Step 1: parameters.} $\lambda$ is a Harish-Chandra parameter (\lref{Formalism.MainNotions.StandardLattice.line_param_harish_chandra}),
and $p,q\ge1$ (\lref{Formalism.MainNotions.StandardLattice.signature_counts_ge_one}). From $e^{-N}\le\eps_0<\frac12$
and $h_0\ge8192\min(p,q)(pq)^3\log N/\eps_0^2$ we get $h_0>32Npq+2$: indeed $N\ge2$, $\log N>0.69$, $\min(p,q)\ge1$, $pq\ge1$
and $\eps_0^2<\frac14$, so $h_0\ge4\cdot8192\cdot0.69\,(pq)^3>32Npq+2$ because $N\le2pq$. In particular
$h_0>\max(2N,8N\min(p,q))$.

\emph{Step 2: a discrete series realising $\shcurvDs(\rho)$.} By (Existence) there is a discrete series representation
$\sigma$ with parameter $\lambda$, with some formal degree $d>0$ with respect to $\mu_\eps$; then $d/s$ is a formal degree
with respect to $s\mu_\eps$, and the formal degree tie gives $d/s=\shcurvDs(\rho)$. Choose $w\ne0$ in the space of $\sigma$. By Schur
orthogonality, $\|w\|^4/\int|\langle\sigma(g)w,w\rangle|^2d(s\mu_\eps)=\shcurvDs(\rho)$.

\emph{Step 3: the bounds.} By \cref{pf:sh-curv:hc}, $\shcurvDs(\rho)=\prod_{a<b}|\lambda_a-\lambda_b|/(b-a)$, and this equals
$\dim V(\lambda)$ because $\lambda$ is strictly decreasing
(\lref{Formalism.MainNotions.StandardLattice.dimV_eq_fdegHC}). As $\dim V(\lambda)\ge0$ and
$\eps_0>0$, the lower bound holds. For the upper bound, \cref{pf:sh-curv:fdegle} gives
$\shcurvDs(\rho)\le\shcurvDt(\eps,\lambda)|\lambda|^{pq}$; since $|\lambda|\ge1$, $3pq\ge1$ and $\shcurvDt(\eps,\lambda)>0$
(\cref{pf:sh-curv:fdegle}),
$\shcurvDt|\lambda|^{pq}\le\shcurvDt(3pq|\lambda|)^{pq}\le4\shcurvDt(3pq|\lambda|)^{pq}$.

\emph{Step 4: where the minimal $K$-type enters.} The formal proof
states the upper bound first with $\dim S_\Lambda$ in place of $\shcurvDt(\eps,\lambda)$, where $S_\Lambda$ is the sum of all
$K_\eps$-types of highest weight $\Lambda$ in $\sigma$, and identifies $\dim S_\Lambda=\shcurvDt(\eps,\lambda)$: since $h_0>Npq+1$, (Minimal $K$-type) shows
that $S_\Lambda$ is the unique $K_\eps$-type of highest weight $\Lambda$, and (Weyl), applied to the restriction of $\sigma$
to $K_\eps$ and the weight $\Lambda$ with $\Lambda+\rho_c=\lambda+\rho_n$, gives its dimension. This is the only place where
these two inputs are used.
\end{proof}

\begin{remark}\label{pf:sh-curv:dsig-rem}
(1) The proof shows the equality $\shcurvDs(\rho)=\dim V(\lambda(\rho))$ in large gap; the lower bound with the factor
$1-\eps_0$ and the factors $4$ and $3pq$ in the upper bound are therefore not sharp, and the hypotheses $e^{-N}\le\eps_0$
and the precise form of $\shcurvGap$ are only used to guarantee $\shcurvGap(\eps_0)>32Npq+2$. (2) The eventual
formulation over families is not needed: the threshold is $0$.
\end{remark}

\subsubsection{The slope of the ramified lattice}

\begin{definition}[The ramified lattice]\label{pf:sh-curv:ramdef}
Let $D$ be a datum and $u$ a word of $D$. We use the hermitian lattices of \cref{sec:pf:top-shape}: $M_u$ is the hermitian
lattice at level $0$ attached to $u$, over the ring of integers of the coefficient CM field $\mathit{LL}$, with its
metrics at the infinite places of $\mathit{LL}$ and its Arakelov slope $\mu(M_u)$. Let $T=T_u$ be the level prime of $u$.
The datum provides a set $\mathcal V_T$ of $f$ places with residue fields $\F_v$ of $T$ elements, and actions
$a\mapsto\rho_T(v,a)$ of $\F_v^\times$ (\emph{diamond operators}) on the space of functions on $\mathrm{cls}_0\times Y_0$,
where $\mathrm{cls}_0$ is the set of classes at level $0$; an embedding $\mathrm{emb}_0$ identifies the space $V_u$ of
$M_u$ with a space of such functions contained in the space $M_{u,0}$ of forms at level $0$. Put
$\Nm_v=\sum_{a\in\F_v^\times}\rho_T(v,a)$ and let the \emph{ramified subspace} be
$V_u^{\shcurvRam}=\mathrm{emb}_0^{-1}\bigl(M_{u,0}\cap\bigcap_v\ker\Nm_v\bigr)$. The \emph{ramified lattice}
$M_u^{\shcurvRam}$ is $M_u\cap V_u^{\shcurvRam}$ with the metrics of $M_u$.
\formal{Formalism.BasicNotions.ColumnsAtT.Mram}
\formal{Formalism.BasicNotions.MemberForms.ramLat}
\end{definition}

\begin{theorem}[The slope line]\label{pf:sh-curv:slopeline}
Let $D$ be a datum and $u$ a word such that
\begin{enumerate}[label=(\roman*)]
\item the Hida setting conditions hold at $u$: $T>2fN(m)+1$, $4a_{\max}(u)<T$, the alcove condition and $b_{\mathrm{top}}(\tau)
\ge0$ for all tops $\tau$;
\item the concrete forms at level zero exist: there is a tied symmetric system of forms at $u$ whose level-zero character
$\chi$ has finite order and takes values in $\mathit{LL}$, and the stabilisers at level zero are trivial.
\end{enumerate}
Then
\[
\mu\bigl(M_u^{\shcurvRam}\bigr)=\mu(M_u)-\frac{f\,\log(T-1)}{T-2}.
\]
\formal{Formalism.MainNotions.RT.SlopeLineT}
\formal{Formalism.MainNotions.product_first_zeroth_parts.SlopeLine}
\formallabel{Formalism.MainNotions.slopeLineT_holds}
\formallabel{Formalism.L2.GlobalTheoryTheory.LocalLatticeBases.slopeLine_0}
\formallabel{Formalism.L2.GlobalTheoryTheory.LocalLatticeBases.S0.slopeLine_tied}
\formallabel{Formalism.L2.GlobalTheoryTheory.LocalLatticeBases.S0.slopeLine_of_tied}
\formallabel{Formalism.L2.GlobalTheoryTheory.LocalLatticeBases.S0.slope_formula_of_character_finite}
\formallabel{Formalism.L2.GlobalTheoryTheory.LocalLatticeBases.S0.slope_formula_of_isometry}
\formallabel{Formalism.L2.GlobalTheoryTheory.LocalLatticeBases.S0.slope_formula_of_split}
\formallabel{Formalism.L2.GlobalTheoryTheory.LocalLatticeBases.S0.slope_formula_of_diamonds}
\formallabel{Formalism.L2.GlobalTheoryTheory.LocalLatticeBases.S0.slope_formula_of_compatible_slice}
\formallabel{Formalism.L2.GlobalTheoryTheory.LocalLatticeBases.S0.slope_formula_of_slice}
\formallabel{Formalism.L2.GlobalTheoryTheory.LocalLatticeBases.S0.slope_formula_of_induced}
\formallabel{Formalism.L2.GlobalTheoryTheory.LocalLatticeBases.S0.slice_summands_orthogonal}
\formallabel{Formalism.L2.GlobalTheoryTheory.LocalLatticeBases.S0.group_rep_majorant_isometry}
\formallabel{Formalism.L2.GlobalTheoryTheory.LocalLatticeBases.S0.ramified_sub_eq_inf_norm_ker}
\formallabel{Formalism.L2.GlobalTheoryTheory.LocalLatticeBases.S0.three_le_level_prime}
\end{theorem}

The formula is a statement about induced lattices. Let $G=\prod_{i\in I}G_i$ be a product of finite abelian groups and $W$
a hermitian lattice over $\OO_K$ ($K$ a CM field) with space $W_K$. Functions $F\colon G\to W_K$ form the induced space;
for $i\in I$ let $\nu_i$ be the map summing $F$ over the cosets of $G_i$, and let
$\mathrm{Ram}(G,W_K)=\bigcap_i\ker\nu_i$ (the functions all of whose partial sums along the factors vanish).

\begin{lemma}[Slopes of ramified parts of induced lattices]\label{pf:sh-curv:induced}
Let $B$ be a hermitian lattice (with a good frame), $W$ a non-zero hermitian lattice, $q\ge3$, $G=\prod_{i\in I}G_i$ with
$|G_i|=q-1$ for all $i$, and $\Psi$ an injective linear map from the functions $G\to W_K$ to the space of $B$ such that
\begin{enumerate}[label=(\alph*)]
\item $\Psi$ is isometric for the sum metric: $\langle\Psi F,\Psi F'\rangle_v=\sum_{g}\langle F(g),F'(g)\rangle_v$ at every
infinite place $v$;
\item $B=\Psi\bigl(\{F:F(g)\in W\ \forall g\}\bigr)$.
\end{enumerate}
Then the sublattice $B^{\shcurvRam}=B\cap\Psi(\mathrm{Ram}(G,W_K))$ satisfies
$\mu(B^{\shcurvRam})=\mu(B)-|I|\log(q-1)/(q-2)$.
\formallabel{Formalism.L2.GlobalTheoryTheory.LocalLatticeBases.S0.slope_ramified_induced_prime_field}
\formallabel{Formalism.L2.GlobalTheoryTheory.LocalLatticeBases.S0.slope_ramified_induced}
\formallabel{Formalism.L2.GlobalTheoryTheory.LocalLatticeBases.S0.slope_eq_of_isometric_sum}
\end{lemma}

\begin{proof}
We use the \emph{Kronecker slope formula} (\lref{Formalism.L2.Col.T2.kron_mu}): if $\beta$ is a non-empty finite set,
$\Phi$ an injective linear map from $W_K^\beta$ to the space of $B$, and $g\in M_\beta(\Z)$ a matrix with $\det g>0$ such
that $\langle\Phi F,\Phi F'\rangle_v=\sum_{b,c}g_{bc}\langle F_b,F'_c\rangle_v$ at all infinite places, then the image
$\Phi(W^\beta)$ (with the metrics of $B$) has slope $\mu(W)-\log(\det g)/|\beta|$.

\emph{The slope of $B$.} Applied to $\Phi=\Psi$ and $g=1$, this gives $\mu(B)=\mu(W)$, by (a) and (b).

\emph{A basis of the ramified functions.} For each $i$ fix $1\in G_i$ and put $G_i'=G_i\smallsetminus\{1\}$ and
$\beta=\prod_iG_i'$. For $c=(c_i)\in\beta$ let $s_c\in\Z^G$ be $\bigotimes_i(\delta_{c_i}-\delta_1)$, i.e.\
$s_c(g)=\prod_i(\delta_{g_i,c_i}-\delta_{g_i,1})$, and let $S\in M_{\beta\times G}(\Z)$ be the matrix with rows $s_c$. The
vectors $s_c$ form a $\Z$-basis of the functions $G\to\Z$ all of whose partial sums along the factors vanish (a
tensor product of the bases $\delta_c-\delta_1$, $c\in G_i'$, of the sum-zero functions on the $G_i$). Consequently the map
$\kappa_S\colon W_K^\beta\to W_K^G$, $(F_c)_c\mapsto\sum_cs_c\otimes F_c$, is injective with image $\mathrm{Ram}(G,W_K)$,
and an element of its image has all values in $W$ if and only if all $F_c\in W$
(\lref{Formalism.L2.GlobalTheoryTheory.LocalLatticeBases.S0.induced_lattice_map_sz}). Hence
$B^{\shcurvRam}=(\Psi\circ\kappa_S)(W^\beta)$.

\emph{The Gram determinant.} By (a), $\Psi\circ\kappa_S$ satisfies the hypothesis of the Kronecker slope formula with
$g=SS^{\mathsf T}=\bigotimes_i(1+J)$, where $1+J$ is the $(q-2)\times(q-2)$ matrix with $2$ on the diagonal and $1$
elsewhere. Since $\det(1+J)=q-1$ (the eigenvalues of $J$ are $q-2$ and $0$) and
$\det(A_1\otimes A_2)=\det(A_1)^{n_2}\det(A_2)^{n_1}$ for square matrices of sizes $n_1,n_2$,
\[
\log\det g=|I|\,(q-2)^{|I|-1}\log(q-1),\qquad|\beta|=(q-2)^{|I|},
\]
so $\log(\det g)/|\beta|=|I|\log(q-1)/(q-2)$ (\lref{Formalism.L2.Col.T2.log_det_szGram_div}). The Kronecker slope
formula gives $\mu(B^{\shcurvRam})=\mu(W)-|I|\log(q-1)/(q-2)$, and $\mu(W)=\mu(B)$.
\end{proof}

\begin{proof}[Proof of \cref{pf:sh-curv:slopeline}]
Write $B=M_u$, with space $V=V_u$. We first record the inputs.
\begin{enumerate}[label=(\alph*)]
\item From the tied forms, the stabilisers at level zero are trivial
(\lref{Formalism.L2.GlobalTheoryTheory.LocalLatticeBases.S0.admissible_stabilisers_of_compatible_forms}); the
corresponding part of hypothesis (ii) is not used.
\item The level prime splits as required in the tied forms
(\lref{Formalism.L2.GlobalTheoryTheory.LocalLatticeBases.S0.splitAtT_of_tied}), and consequently there are
\emph{diamond elements}: a homomorphism $\rho$ from $G=\prod_{v\in\mathcal V_T}\F_v^\times$ to $\End_{\mathit{LL}}(V)$ with
$\mathrm{emb}_0(\rho(a_v)x)=\rho_T(v,a)(\mathrm{emb}_0x)$, where $a_v\in G$ is $a$ in the factor $v$
(\lref{Formalism.L2.GlobalTheoryTheory.LocalLatticeBases.S0.diamond_elements_of_split}). The group $G$ also acts
on $\mathrm{cls}_0$, compatibly: $\mathrm{emb}_0(\rho(g)x)(c,\cdot)=e_{g,c}\bigl(\mathrm{emb}_0(x)(g\cdot c,\cdot)\bigr)$ for an
invertible linear map $e_{g,c}$ of the functions on $Y_0$
(\lref{Formalism.L2.GlobalTheoryTheory.LocalLatticeBases.S0.group_rep_transport_value}), and the action is free
because the stabilisers are trivial (\lref{Formalism.L2.GlobalTheoryTheory.LocalLatticeBases.S0.group_action_injective}).
\item (Diamond isometry) The operators $\rho(a_v)$ preserve the metrics of $B$ at all infinite places
(\lref{Formalism.L2.GlobalTheoryTheory.LocalLatticeBases.S0.diamond_isometry_at_infinity}; this uses that $\chi$ has
finite order and values in $\mathit{LL}$). By multiplicativity, every $\rho(g)$ does.
\item (Disjoint supports) If $\mathrm{emb}_0x$ and $\mathrm{emb}_0y$ have disjoint supports in $\mathrm{cls}_0$, then $x$ and
$y$ are orthogonal at every infinite place
(\lref{Formalism.L2.GlobalTheoryTheory.LocalLatticeBases.S0.metric_disjoint_support_orthogonal}).
\end{enumerate}
\emph{The slice.} Choose a representative of each $G$-orbit on $\mathrm{cls}_0$ and let $W_0\subset V$ be the subspace of
vectors $x$ with $\mathrm{emb}_0x$ supported on the representatives. Let $\Psi(F)=\sum_g\rho(g)F(g)$ for $F\colon G\to
W_0$. By (b), $\mathrm{emb}_0(\rho(g)x)$ at a class $c$ is the image of $\mathrm{emb}_0(x)$ at $g\cdot c$ under an
invertible linear map of the functions on $Y_0$; since the action is free, $\mathrm{emb}_0(\Psi F)$ at $c$ therefore only
involves the summand indexed by the unique $g$ with $g\cdot c$ a representative. Hence:
$\Psi$ is bijective (\lref{Formalism.L2.GlobalTheoryTheory.LocalLatticeBases.S0.slice_map_bijective}; injectivity
follows from the previous remark and the injectivity of $\mathrm{emb}_0$, surjectivity from the fact that the restriction of
a form at level zero to a set of classes is again the image of an element of $V$, which uses the triviality of the
stabilisers); $\langle\rho(g)x,\rho(h)y\rangle_v=\delta_{g,h}\langle x,y\rangle_v$ for $x,y\in W_0$ (by (c) for $g=h$, and
by (d) and freeness for $g\ne h$, since $\rho(g)x$ and $\rho(h)y$ are supported on disjoint translates of the
representatives); $\Psi(F)\in B$ if and only if $F(g)\in B$ for all $g$
(\lref{Formalism.L2.GlobalTheoryTheory.LocalLatticeBases.S0.slice_mem_lattice_iff}: membership in $B$ is a local
integrality condition on the values of $\mathrm{emb}_0$ at each class, and the $\rho(g)$ preserve $B$); and $B\ne0$
(\lref{Formalism.L2.GlobalTheoryTheory.LocalLatticeBases.S0.lattice_level_zero_ne_bot}).

\emph{Induction.} From these properties one obtains a non-zero hermitian lattice $W$ and an injective map $\Psi'$ from
functions $G\to W_K$ to $V$ satisfying (a), (b) of \cref{pf:sh-curv:induced} and
$\Psi'(\mathrm{Ram}(G,W_K))=\bigcap_v\ker(\sum_{a}\rho(a_v))$
(\lref{Formalism.L2.GlobalTheoryTheory.LocalLatticeBases.S0.induced_of_slice}, which builds $W$ from $W_0$ and $B$). Moreover
$V^{\shcurvRam}=\bigcap_v\ker(\sum_a\rho(a_v))$, because $\mathrm{emb}_0$ is injective with image in $M_{u,0}$ and
intertwines $\sum_a\rho(a_v)$ with $\Nm_v$. So
$M_u^{\shcurvRam}=B\cap\Psi'(\mathrm{Ram}(G,W_K))$.

\emph{Conclusion.} The Hida setting gives $T>2fN(m)+1\ge3$, as $f\ge2$ and $N(m)\ge1$; moreover $|\F_v^\times|=T-1$ and
$|\mathcal V_T|=f$. \Cref{pf:sh-curv:induced} with $q=T$ and $I=\mathcal V_T$ gives the formula.
\end{proof}

\begin{remark}\label{pf:sh-curv:slopeline-rem}
Of the Hida setting conditions (i) only $T\ge3$ is used, and of (ii) only the tied forms and the finiteness of $\chi$;
the triviality of the stabilisers assumed in (ii) is re-derived from the forms.
\end{remark}

% ===========================================================================
\subsection{The extremal projector and the tool statements}\label{sec:pf:sh-curv:ep}

The results of this subsection are the inputs $(\mathrm E)$, the package of tool statements, and the threshold functions
of the first theorem, used in the proof of the local lattice theorem in \cref{sec:pf:linkG}
(\cref{pf:linkG:hyps,pf:linkG:locmodel}).

\subsubsection{The extremal projector}

\begin{definition}[Truncated extremal projectors]\label{pf:sh-curv:epdef}
Let $K$ be a field of characteristic $0$, $n\ge1$ and $V$ a finite-dimensional $K$-vector space. A
\emph{$\mathfrak{gl}_n$-representation} on $V$ is a family $(E_{ab})_{a,b<n}$ of endomorphisms of $V$ with
$[E_{ab},E_{cd}]=\delta_{bc}E_{ad}-\delta_{da}E_{cb}$, such that every vector lies in a finite-dimensional subspace
stable under all $E_{ab}$ and $V$ is the sum of its \emph{weight spaces}
$V_\mu=\{x:E_{aa}x=\mu_ax\ \forall a\}$, $\mu\in\Z^n$. Write $\shcurvNm V=\sum_{a>b}E_{ab}V$ for the span of the images of
the lowering operators, $E^{(k)}_{ab}=E_{ab}^k/k!$, and $\alpha_{ab}=e_a-e_b\in\Z^n$. A \emph{normal ordering} is a list
$\mathrm{ord}$ of the pairs $(a,b)$ with $a<b$, each occurring once, such that for $a<b<c$ the pair $(a,c)$ lies between
$(a,b)$ and $(b,c)$. For $m,T\in\N$, a list $\mathrm{ord}$ of pairs and $\mu\in\Z^n$ put
\begin{equation}\label{eq:pf:sh-curv:pfac}
p_{ab}(m,T)=\sum_{k=0}^T\frac{(-1)^k}{\binom{m+k}{k}}\,E^{(k)}_{ba}E^{(k)}_{ab},\qquad
\shcurvPP_\mu=\shcurvPP(\mu,\mathrm{ord},T)=\prod_{(a,b)\in\mathrm{ord}}p_{ab}(m_{ab},T),
\end{equation}
where $m_{ab}=\max(0,\mu_a-\mu_b+b-a)$, the product taken in the order of the list (the first pair is leftmost).
\formal{Formalism.MainNotions.G.IsGLRep}
\formal{Formalism.MainNotions.G.weightSp}
\formal{Formalism.MainNotions.G.nMinus}
\formal{Formalism.MainNotions.G.NormalOrder}
\formal{Formalism.MainNotions.G.ppFac}
\formal{Formalism.MainNotions.G.pp}
\end{definition}

If $\mu_a-\mu_b+b-a\ge1$ for $a<b$, then on $V_\mu$ the factor $p_{ab}(m_{ab},T)$ is the truncation at $T$ of the series
$\sum_k(-1)^k\bigl(k!\prod_{j=1}^k(h_{ab}+\rho(h_{ab})+j)\bigr)^{-1}f^k_{ab}e^k_{ab}$ of Asherova, Smirnov and Tolstoy, with
$e_{ab}=E_{ab}$, $f_{ab}=E_{ba}$, $h_{ab}=E_{aa}-E_{bb}$ and $\rho(h_{ab})=b-a$, and $\shcurvPP_\mu$ is the truncated
extremal projector \cite{AsherovaSmirnovTolstoy1971,Zhelobenko1989}.

\begin{theorem}[The extremal projector statement $(\mathrm E)$]\label{pf:sh-curv:E}
Let $K$ be a field of characteristic $0$, $V$ a finite-dimensional $\mathfrak{gl}_n$-representation over $K$, $\mu\in\Z^n$
with $\mu_a-\mu_b+b-a\ge1$ for all $a<b$, $\mathrm{ord}$ a normal ordering and $T\ge\dim V$. Then
\begin{enumerate}[label=(\roman*)]
\item $\shcurvPP_\mu V_\mu\subseteq V_\mu$;
\item $\shcurvPP_\mu x=0$ for $x\in V_\mu\cap\shcurvNm V$;
\item $\shcurvPP_\mu x=x$ for $x\in V_\mu$ with $E_{ab}x=0$ for all $a<b$.
\end{enumerate}
(Formally the statement is indexed by a datum $D$, which does not occur in it.)
\formal{Formalism.MainNotions.CovariantDegrees.adm5}
\formallabel{Formalism.MainNotions.G.ExtremalProjector.extremal_projector_statement}
\formallabel{Formalism.MainNotions.G.GLWeightTheory.sizes_bound_of_log_estimates}
\end{theorem}

This is the classical theorem on the extremal projector of $\mathfrak{gl}_n$
\cite{AsherovaSmirnovTolstoy1971}, \cite[Thms.~2 and 4]{Zhelobenko1989}: the (untruncated) projector does not depend on the
normal ordering, satisfies $e_{ab}\shcurvPP=0$ for all $a<b$, and projects onto the highest weight vectors along
$\shcurvNm V$. The formal statement differs only in that the series are truncated at $T\ge\dim V$ (which does not change
them on $V$, see \cref{pf:sh-curv:nil}) and the weight is fixed in advance. We give the proof of the formal development; the
two combinatorial inputs (\cref{pf:sh-curv:conn,pf:sh-curv:yb}) are cited.

\begin{lemma}[Nilpotency]\label{pf:sh-curv:nil}
For $a\ne b$ and $k\ge\dim V$, $E_{ab}^k=0$.
\formallabel{Formalism.MainNotions.G.ExtremalProjector.restNil}
\end{lemma}

\begin{proof}
$E_{ab}$ maps $V_\nu$ to $V_{\nu+\alpha_{ab}}$ (from the commutation relation with $E_{cc}$), and $V$ has finitely many
weights; hence $E_{ab}^kV_\nu=0$ for $k$ large, $E_{ab}$ is nilpotent, and a nilpotent endomorphism of a space of dimension
$d$ satisfies $E^d=0$.
\end{proof}

\begin{lemma}[Elementary properties]\label{pf:sh-curv:elem}
Let $\mu,\nu\in\Z^n$, $\mathrm{ord}$ a list of pairs $(a,b)$ with $a<b$ and $T\in\N$.
\begin{enumerate}[label=(\alph*)]
\item $\shcurvPP(\mu,\mathrm{ord},T)V_\nu\subseteq V_\nu$.
\item If $E_{ab}x=0$ for all $(a,b)\in\mathrm{ord}$, then $\shcurvPP(\mu,\mathrm{ord},T)x=x$.
\item $\shcurvPP(\mu,\mathrm{ord},T)x-x\in\shcurvNm V$ for all $x$.
\item If $y$ satisfies $E_{a,a+1}y=0$ for all $a<n-1$, then $E_{ab}y=0$ for all $a<b$.
\end{enumerate}
\formallabel{Formalism.MainNotions.G.GLWeightTheory.extremal_projector_preserves_weight}
\formallabel{Formalism.MainNotions.G.GLWeightTheory.extremal_projector_fixes_hwv}
\formallabel{Formalism.MainNotions.G.GLWeightTheory.extremal_projector_sub_mem_nMinus}
\formallabel{Formalism.MainNotions.G.GLWeightTheory.simple_raising_kill_all_raising}
\end{lemma}

\begin{proof}
(a) $E^{(k)}_{ba}E^{(k)}_{ab}$ maps $V_\nu$ to $V_{\nu+k\alpha_{ab}-k\alpha_{ab}}=V_\nu$, so every factor preserves $V_\nu$.
(b) If $E_{ab}x=0$, all terms with $k\ge1$ of $p_{ab}x$ vanish and $p_{ab}x=x$; induct along the list. (c) For $a<b$, the
terms with $k\ge1$ of $p_{ab}x$ lie in $E_{ba}V\subseteq\shcurvNm V$, so $p_{ab}x-x\in\shcurvNm V$; then
$\shcurvPP x-x$ is a sum of such differences. (d) By induction on $b-a$: if $b=c+1$ with $a<c$, then $E_{ac}y=0$ and
$E_{cb}y=0$, and $E_{ab}y=E_{ac}E_{cb}y-E_{cb}E_{ac}y=0$ by the commutation relation.
\end{proof}

\begin{lemma}\label{pf:sh-curv:hwvn}
If $y\in V_\mu\cap\shcurvNm V$ and $E_{ab}y=0$ for all $a<b$, then $y=0$.
\formallabel{Formalism.MainNotions.G.GLWeightTheory.hwv_mem_nMinus_eq_zero}
\end{lemma}

\begin{proof}
Suppose $y\ne0$. Let $X\subseteq V$ be the span of the vectors $uv$, $u$ a product of the $E_{ab}$ and $v$ a highest weight
vector of weight $\mu$ (the \emph{isotypic subspace} of $\mu$); it is stable and contains $y$, and its weights are
$\le\mu$ in the dominance order. By complete reducibility (\lref{Formalism.MainNotions.G.GLWeightTheory.cr_compl}:
every stable subspace of a finite-dimensional representation which is the sum of its weight spaces has a stable
complement; Weyl's theorem) there is a stable $X'$ with $V=X\oplus X'$. Writing
$w=p+q$ with $p\in X$, $q\in X'$, every lowering operator maps $w$ to $E_{ab}p+E_{ab}q$ with $E_{ab}p\in X$ of weights
$<\mu$ and $E_{ab}q\in X'$; hence $y=s+x'$ with $s\in X$ a sum of vectors of weights $<\mu$ and $x'\in X'$. Since $y\in X$,
$x'=y-s\in X\cap X'=0$, so the non-zero vector $y$ of weight $\mu$ is a sum of vectors of weights $\ne\mu$, which is absurd.
\end{proof}

The heart of the proof is the vanishing of the simple raising operators on the image of $\shcurvPP_\mu$. Write
$p_{ab}(m)=p_{ab}(m,T)$ for a fixed $T\ge\dim V$. Two lists of positive roots are related by a \emph{move} if one is
obtained from the other by exchanging two adjacent pairs $(a,b),(c,d)$ with $\{a,b\}\cap\{c,d\}=\emptyset$, or by replacing
three adjacent pairs $(a,b),(a,c),(b,c)$ with $a<b<c$ by $(b,c),(a,c),(a,b)$, or conversely.

\begin{lemma}[Connectedness of normal orderings]\label{pf:sh-curv:conn}
For every normal ordering $\mathrm{ord}$ and every $a$ with $a+1<n$ there is a list $l$ such that $((a,a+1))\cdot l$ is
obtained from $\mathrm{ord}$ by a finite sequence of moves.
\formallabel{Formalism.MainNotions.G.ExtremalProjector.restConn}
\end{lemma}

\begin{proof}
Normal orderings of the positive roots of $\GL_n$ correspond to reduced decompositions of the longest element of the
symmetric group \cite{Zhelobenko1989}, and the moves correspond to the commutation and braid relations; the statement is
then Matsumoto's theorem, since the longest element has a reduced decomposition beginning with any simple reflection. The
formal proof is a direct sorting argument on normal orderings
(\lref{Formalism.MainNotions.G.ExtremalProjector.cn_sort} and \lref{Formalism.MainNotions.G.ExtremalProjector.cn_iter}).
\end{proof}

\begin{lemma}[Commuting factors]\label{pf:sh-curv:comm}
If $\{a,b\}\cap\{c,d\}=\emptyset$, then $p_{ab}(s)$ and $p_{cd}(s')$ commute for all $s,s'$.
\formallabel{Formalism.MainNotions.G.ExtremalProjector.restComm}
\end{lemma}

\begin{proof}
By the commutation relations, $E_{ab}$ and $E_{ba}$ commute with $E_{cd}$ and $E_{dc}$ when the four indices satisfy
$a\ne c,d$ and $b\ne c,d$; hence so do the polynomials $p_{ab}(s)$ and $p_{cd}(s')$ in them.
\end{proof}

\begin{lemma}[The braid relation]\label{pf:sh-curv:yb}
Let $a,b,c$ be distinct, $\nu\in\Z^n$, and $s,t\in\N$ with $s=\nu_a-\nu_b+\sigma+1$ and $t=\nu_b-\nu_c+\tau+1$ for some
$\sigma,\tau\in\N$. Then on $V_\nu$
\[
p_{ab}(s)\,p_{ac}(s+t)\,p_{bc}(t)=p_{bc}(t)\,p_{ac}(s+t)\,p_{ab}(s).
\]
\formallabel{Formalism.MainNotions.G.ExtremalProjector.restYB}
\formallabel{Formalism.MainNotions.G.ExtremalProjector.restYB_of}
\end{lemma}

\begin{proof}
This is the factorisation identity for the extremal projector of $\mathfrak{gl}_3$
\cite{AsherovaSmirnovTolstoy1971,Zhelobenko1989}. The formal proof proceeds as follows. The case $\sigma=\tau=0$
(\lref{Formalism.MainNotions.G.ExtremalProjector.ybat_base}) is deduced from \cref{pf:sh-curv:kill} and two transport
identities (\lref{Formalism.MainNotions.G.ExtremalProjector.restPair}): for $w\in V_\nu$,
$E_{bc}\,p_{ab}(s)p_{ac}(u)w=p_{ab}(s)p_{ac}(u)E_{bc}w$ if $u-(\nu_a-\nu_c)-1=s-(\nu_a-\nu_b)$, and
$E_{ab}\,p_{bc}(t)p_{ac}(u)w=p_{bc}(t)p_{ac}(u)E_{ab}w$ if $u-(\nu_a-\nu_c)-1=t-(\nu_b-\nu_c)$. The passage from $(\sigma,\tau)$ to
$(\sigma+1,\tau)$ (\lref{Formalism.MainNotions.G.ExtremalProjector.ybat_step}) applies the identity for $(\sigma,\tau)$ to
the representation $V\otimes K^n$ ($K^n$ the standard representation) and the vector $x\otimes e_a$, of weight
$\nu+e_a$, and reads off the component along $e_a$; the factors $p_{ab},p_{ac},p_{bc}$ act on this component as on $V$. A
symmetry exchanging $\sigma$ and $\tau$ (\lref{Formalism.MainNotions.G.ExtremalProjector.ybat_symm}) then gives all
$(\sigma,\tau)$ by induction.
\end{proof}

\begin{lemma}[The $\mathfrak{sl}_2$ computation]\label{pf:sh-curv:kill}
Let $a\ne b$, $\nu\in\Z^n$ and $s\in\N$ with $s=\nu_a-\nu_b+1$. Then $E_{ab}\,p_{ab}(s)x=0$ for all $x\in V_\nu$.
\formallabel{Formalism.MainNotions.G.ExtremalProjector.restKill}
\end{lemma}

\begin{proof}
Put $e=E_{ab}$, $f=E_{ba}$, $h=E_{aa}-E_{bb}$, so $[e,f]=h$, $[h,f]=-2f$, and $e f^{(k)}=f^{(k)}e+f^{(k-1)}(h-k+1)$ for $k\ge1$
(induction on $k$). The vector $e^{(k)}x$ has $h$-eigenvalue $s-1+2k$, and $e\,e^{(k)}=(k+1)e^{(k+1)}$. Hence
\[
e\,f^{(k)}e^{(k)}x=(k+1)f^{(k)}e^{(k+1)}x+(s+k)f^{(k-1)}e^{(k)}x\qquad(f^{(-1)}=0).
\]
With $c_k=(-1)^k/\binom{s+k}{k}$ we have $c_{k+1}(s+k+1)=-(k+1)c_k$, so in $e\,p_{ab}(s)x=\sum_{k\le T}c_k\,e
f^{(k)}e^{(k)}x$ the coefficient of $f^{(k)}e^{(k+1)}x$ is $(k+1)c_k+(s+k+1)c_{k+1}=0$ for $k<T$, and what remains is
$(T+1)c_Tf^{(T)}e^{(T+1)}x$, which vanishes by \cref{pf:sh-curv:nil} since $T+1>\dim V$.
\end{proof}

\begin{proposition}[Simple raising operators kill the image]\label{pf:sh-curv:simplekill}
In the situation of \cref{pf:sh-curv:E}, $E_{a,a+1}\shcurvPP_\mu x=0$ for all $x\in V_\mu$ and $a+1<n$.
\formallabel{Formalism.MainNotions.G.ExtremalProjector.raising_operators_kill_image}
\formallabel{Formalism.MainNotions.G.ExtremalProjector.raising_kill_of_interface}
\formallabel{Formalism.MainNotions.G.ExtremalProjector.pp_step}
\end{proposition}

\begin{proof}
Write $\shcurvPP(l)=\prod_{(a,b)\in l}p_{ab}(m_{ab})$ for a list $l$, with $m_{ab}=\mu_a-\mu_b+b-a$. We claim that
$\shcurvPP(l)x=\shcurvPP(l')x$ for $x\in V_\mu$ whenever $l'$ is obtained from $l$ by a move. For the exchange of two pairs
with disjoint indices this is \cref{pf:sh-curv:comm}. For the replacement of $(a,b),(a,c),(b,c)$ by $(b,c),(a,c),(a,b)$
(or conversely) inside $l=l_1\cdot((a,b),(a,c),(b,c))\cdot l_2$: the vector $w=\shcurvPP(l_2)x$ lies in $V_\mu$
(\cref{pf:sh-curv:elem}(a)), and $m_{ac}=m_{ab}+m_{bc}$, with $m_{ab}=\mu_a-\mu_b+(b-a-1)+1$ and
$m_{bc}=\mu_b-\mu_c+(c-b-1)+1$; so \cref{pf:sh-curv:yb} with $\sigma=b-a-1$, $\tau=c-b-1$ gives the claim. By
\cref{pf:sh-curv:conn}, $\shcurvPP(\mathrm{ord})x=\shcurvPP((a,a+1)\cdot l)x=p_{a,a+1}(m_{a,a+1})\shcurvPP(l)x$ for some list
$l$. The vector $y=\shcurvPP(l)x$ lies in $V_\mu$ and $m_{a,a+1}=\mu_a-\mu_{a+1}+1$, so $E_{a,a+1}p_{a,a+1}(m_{a,a+1})y=0$ by
\cref{pf:sh-curv:kill}.
\end{proof}

\begin{proof}[Proof of \cref{pf:sh-curv:E}]
(i) is \cref{pf:sh-curv:elem}(a), and (iii) is \cref{pf:sh-curv:elem}(b), since a normal ordering consists of pairs $(a,b)$
with $a<b$. For (ii), let $x\in V_\mu\cap\shcurvNm V$ and $y=\shcurvPP_\mu x$. Then $y\in V_\mu$; by
\cref{pf:sh-curv:simplekill} and \cref{pf:sh-curv:elem}(d), $E_{ab}y=0$ for all $a<b$; and by \cref{pf:sh-curv:elem}(c),
$y=(y-x)+x\in\shcurvNm V$. \Cref{pf:sh-curv:hwvn} gives $y=0$.
\end{proof}

\subsubsection{The tool statements}

\begin{definition}[Integral forms]\label{pf:sh-curv:uzdef}
We use the terminology of \cref{sec:pf:linkG} (\cref{tab:pf:linkG:clauses}): $\OO$ is a discrete valuation ring with
fraction field $K_\OO$ of characteristic $0$ and a uniformiser over a prime $p$; it is \emph{admissible} if $p\ge5$ and
its residue field contains a field with $p^{n'}$ elements for an even $n'\ge2$. For a $\mathfrak{gl}_n$-representation $V$
over $K_\OO$, a \emph{$U_\Z$-lattice} is a finitely generated $\OO$-submodule $M\subseteq V$ stable under the divided powers
$E^{(k)}_{ab}$ ($a\ne b$) and the binomials $\binom{E_{aa}}{k}$; $\ell(M)$ is the largest $k$ such that $\mu-\nu=k\alpha_{ab}$
for two weights $\mu,\nu$ of $K_\OO M$ and some $a\ne b$; $N_0(n)=n(n-1)/2$ and $d_p(n,\ell)=\lfloor\log_p(\ell+n-1)\rfloor$.
For $\lambda\in\Z^n$ the \emph{isotypic subspace} $V[\lambda]$ is spanned by the vectors $uv$, $u$ a product of the $E_{ab}$
and $v$ a highest weight vector of weight $\lambda$; an \emph{isotypic projection} for $\lambda$ is an endomorphism which is
the identity on $V[\lambda]$ and zero on $V[\lambda']$ for $\lambda'\ne\lambda$, and the \emph{invariant projection} is the
isotypic projection for $\lambda=0$. A \emph{weight projection} for $\nu$ is the identity on $V_\nu$ and zero on the
other weight spaces.
\formal{Formalism.MainNotions.G.AdmissibleDVR}
\formal{Formalism.MainNotions.G.UZLattice}
\formal{Formalism.MainNotions.G.ellOf}
\formal{Formalism.MainNotions.G.deF}
\formal{Formalism.MainNotions.G.IsIsoProj}
\formal{Formalism.MainNotions.G.IsWtProj}
\end{definition}

\begin{lemma}[Classical inputs]\label{pf:sh-curv:admitted}
The following hold.
\begin{enumerate}[label=(\alph*)]
\item (Pieri) For $n\ge1$, a non-increasing $\lambda\in\N^n$ and $b\in\N$, with $\rho=(n-1,\dots,1,0)$ and
$a_\alpha=\det(x_i^{\alpha_j})$,
$a_{\lambda+\rho}\,h_b=\sum_\mu a_{\mu+\rho}$, the sum over $\mu\in\N^n$ with $\lambda_j\le\mu_j$,
$\mu_{j+1}\le\lambda_j$ and $|\mu|=|\lambda|+b$; here $h_b$ is the complete homogeneous symmetric polynomial.
\item (Complete reducibility) Every subspace of a finite-dimensional $\mathfrak{gl}_n$-representation over a field of
characteristic $0$ which is stable under all $E_{ab}$ has a stable complement.
\item (Kummer) For a prime $p$, $k\le m$ and $b>\log_pm$, $v_p\binom mk$ is the number of $i\in[1,b)$ with $p^i\le(k\bmod
p^i)+((m-k)\bmod p^i)$.
\item (Gauss) If $A$ is a unique factorisation domain and $w\in A$ is irreducible, then $w\mid FG$ in a polynomial ring
over $A$ implies $w\mid F$ or $w\mid G$.
\end{enumerate}
Together with $(\mathrm E)$ these form the record of inputs $\mathrm{adm}(D)$ of the tool statements.
\formal{Formalism.MainNotions.CovariantDegrees.admitted}
\formallabel{Formalism.MainNotions.G.GLWeightTheory.weight_condition_from_integral_form}
\formallabel{Formalism.MainNotions.G.GLWeightTheory.exponent_condition_holds}
\formallabel{Formalism.MainNotions.G.GLWeightTheory.key_positivity_condition}
\end{lemma}

\begin{proof}
(a) is the Pieri rule for Schur polynomials, multiplied by the Vandermonde determinant $a_\rho$
\cite[I.(5.16)]{Macdonald1995}; in the development it is
\lref{Formalism.MainNotions.G.GLWeightTheory.positivity_condition_holds}, proved by expanding the product of
alternants. (b) is
Weyl's theorem \cite[\S9.3]{FultonHarris1991}; formally it is derived from the weight decomposition
(\lref{Formalism.MainNotions.G.GLWeightTheory.cr_compl}). (c) is Kummer's theorem on carries (the formal proof uses the
corresponding statement of Mathlib). (d): $w$ is prime since $A$ is factorial, so reduction modulo $w$ maps the polynomial
ring onto a polynomial ring over the domain $A/w$, in which $\bar F\bar G=0$ forces $\bar F=0$ or $\bar G=0$.
\end{proof}

\begin{lemma}[The product formula for invariant projections]\label{pf:sh-curv:invformula}
Assume $(\mathrm E)$. Let $K$ be a field of characteristic $0$, $V$ a finite-dimensional $\mathfrak{gl}_n$-representation,
$\mathrm{Inv}$ the invariant projection, $\pi_0$ a weight projection for $0$, $\mathrm{ord}$ a normal ordering and $T\ge\dim V$.
Then $\mathrm{Inv}=\shcurvPP(0,\mathrm{ord},T)\circ\pi_0$.
\formallabel{Formalism.MainNotions.G.GLWeightTheory.invariant_projection_extremal_formula}
\formallabel{Formalism.MainNotions.G.GLWeightTheory.product_formula_of_extremal_projector}
\end{lemma}

\begin{proof}
The weight $\mu=0$ satisfies $\mu_a-\mu_b+b-a\ge1$ for $a<b$, so $(\mathrm E)$ applies to $\shcurvPP_0=\shcurvPP(0,\mathrm{ord},T)$.
By complete reducibility $V$ is the sum of its isotypic subspaces
(\lref{Formalism.MainNotions.G.GLWeightTheory.sum_isotypic_top_of_weight_data}), so it suffices to compare the two sides
on $V[\lambda]$. If $\lambda=0$ and $x\in V[0]$, then $x$ has weight $0$ and is killed by all $E_{ab}$, $a<b$
(\lref{Formalism.MainNotions.G.GLWeightTheory.trivial_isotypic_invariant_holds}: a highest weight vector of weight $0$
spans a trivial subrepresentation); hence $\pi_0x=x$ and
$\shcurvPP_0x=x$ by $(\mathrm E)$(iii), while $\mathrm{Inv}\,x=x$. If $\lambda\ne0$ and $x\in V[\lambda]$, then
$\mathrm{Inv}\,x=0$, and $\pi_0x$ is a vector of weight $0$ in $V[\lambda]$; since $V[\lambda]$ is spanned by lowering words
applied to highest weight vectors of weight $\lambda\ne0$, every such vector lies in $\shcurvNm V$
(\lref{Formalism.MainNotions.G.GLWeightTheory.nontrivial_isotypic_weight_zero_killed}), and $\shcurvPP_0\pi_0x=0$ by
$(\mathrm E)$(ii).
\end{proof}

\begin{theorem}[The tool statements]\label{pf:sh-curv:tools}
Assume $(\mathrm E)$. Then for every datum $D$ (which does not occur in the statements) the following hold.
\begin{enumerate}[label=(T\arabic*)]
\item (Projection theorems) For every admissible $\OO$, every $\mathfrak{gl}_n$-representation $V$ over $K_\OO$ and every
$U_\Z$-lattice $M$ spanning $V$:
\begin{enumerate}[label=(\alph*)]
\item the invariant projection satisfies $\mathrm{Inv}=\shcurvPP(0,\mathrm{ord},T)\circ\pi_0$ for every weight projection
$\pi_0$ for $0$, normal ordering $\mathrm{ord}$ and $T\ge\dim V$, and $p^{N_0(n)d_p(n,\ell(M))}\,\mathrm{Inv}(M)\subseteq M$;
\item for every finite set $S$ of dominant weights and isotypic projections $e_\lambda$ ($\lambda\in S$),
$p^{\max_{\lambda\in S}N_0(n)d_p(n,\ell(M)+\ell(\lambda))}\sum_{\lambda\in S}e_\lambda(M)\subseteq M$, where
$\ell(\lambda)=\lambda_1-\lambda_n$;
\item the comparison of the lattices $\nabla(\lambda,v)\supseteq\Delta(v)$ attached to a highest weight vector, the
translation corollary and the filtered projection corollary (clauses 8, 10 and 11 of \cref{tab:pf:linkG:clauses});
\end{enumerate}
\item (Valuations) the additivity of the valuations $\mathrm{val}_S$ and $\lambda_{1,S}$ of polynomials and their description
by the weights of the generated subrepresentation and of the isotypic components (clause 13);
\item (Determinant divisibility) clause 22 of \cref{tab:pf:linkG:clauses};
\item (Unit products) the expansions of products of $p$-adic unit functions attached to factorials, the factorial unit
decomposition and the degree truncation lemma (clause 24);
\item (Derivatives) the derivative valuation bound, the digit expansion of products of binomial coefficients, and the
stability of digit classes under partial derivatives (clause 25);
\item (Splitting package) seven statements on the polynomials $\Psi$ attached to Gelfand--Tsetlin paths of $\GL_{r+K}$:
the automorphic-side statement, a definite rank-one statement, a rotated sum bound, the path weight components, the
saturation of joint types, the dependence of Witt-lifted substitutions only on the rational plane, and the canonical translate
jet identity;
\item (Path bases) for every discrete valuation ring $\OO$, all $N,n$, letters and block sizes $B$, the polynomials
$\Psi_\OO(P,I)$, for paths $P$ with steps given by the letters and $B$ and standard fillings $I$, are linearly independent
over $\OO$ and span the polynomial space $\mathrm{Pol}_\OO(n,B)$ (clause 15).
\end{enumerate}
\formal{Formalism.MainNotions.CovariantDegrees.Conclusion1}
\formal{Formalism.MainNotions.CovariantDegrees.Conclusion2}
\formal{Formalism.MainNotions.CovariantDegrees.HyperplaneRank}
\formallabel{Formalism.MainNotions.G.tool_statements_hold}
\formallabel{Formalism.MainNotions.G.GLWeightTheory.problem_holds}
\formallabel{Formalism.MainNotions.G.GLWeightTheory.part_one_of_extremal_projector}
\formallabel{Formalism.MainNotions.G.GLWeightTheory.counting_statement_of_block_form}
\formallabel{Formalism.MainNotions.G.GLWeightTheory.invariant_projection_of_product_formula}
\formallabel{Formalism.MainNotions.G.ColumnFamilies.splitting_package_holds}
\formallabel{Formalism.MainNotions.G.GLWeightTheory.true_derivative_digit_class}
\end{theorem}

\begin{proof}
Formally the theorem is $\mathrm{Concl}_1(D)\wedge\mathrm{Concl}_2(D)\wedge\mathrm{HR}(D)$, where $\mathrm{Concl}_1$
collects (T1)--(T5), $\mathrm{Concl}_2$ is (T6) and $\mathrm{HR}$ is (T7). The proof assembles results proved in the
development; we indicate for each component where it comes from and how $(\mathrm E)$ enters.

\emph{(T1)--(T5).} These are derived from $\mathrm{adm}(D)$ (\cref{pf:sh-curv:admitted} and $(\mathrm E)$), of which only
$(\mathrm E)$ is used directly; complete reducibility is used through its own proof.
\begin{itemize}
\item (T1)(a): the product formula is \cref{pf:sh-curv:invformula}; this is the only place where $(\mathrm E)$ is used. The
denominator bound is proved by a separate argument which does not use $(\mathrm E)$: from complete reducibility, the bound
$|\nu_a|\le\ell(M)$ for the weights $\nu$ of $V$ with $\sum_a\nu_a=0$, and a denominator bound for an auxiliary
projector onto the invariants built from divided powers
(\lref{Formalism.MainNotions.G.GLWeightTheory.denominator_lemma_without_projector}).
\item (T1)(b), (c): from complete reducibility and the supremum-norm bound for invariant projections
(\lref{Formalism.MainNotions.G.GLWeightTheory.linf_lemma_without_projector}), by
\lref{Formalism.MainNotions.G.GLWeightTheory.isotypicProjectorTheoremField_of_linfty_bound},
\lref{Formalism.MainNotions.G.GLWeightTheory.integral_form_of_linf_lemma},
\lref{Formalism.MainNotions.G.GLWeightTheory.isotypic_lattice_corollary} and
\lref{Formalism.MainNotions.G.GLWeightTheory.polynomial_isotypic_lattice_corollary}.
\item (T2): \lref{Formalism.MainNotions.G.GLWeightTheory.valuationsAdditive_D7} (additivity) and
\lref{Formalism.MainNotions.G.GLWeightTheory.generated_weight_lemma_of_reducibility} (description by weights, from
complete reducibility).
\item (T3): \lref{Formalism.MainNotions.G.GLWeightTheory.prvE_of_gl_operators_form_representation}, from a height bound
for the weights of the generated representation.
\item (T4): \lref{Formalism.MainNotions.G.GLWeightTheory.mahler_decay_products_of},
\lref{Formalism.MainNotions.G.GLWeightTheory.unit_function_mahler_decay_of} (decay of Mahler coefficients of the unit
functions and of their products) and \lref{Formalism.MainNotions.G.GLWeightTheory.degree_truncation}.
\item (T5): \lref{Formalism.MainNotions.G.GLWeightTheory.derivative_valuation_bound},
\lref{Formalism.MainNotions.G.GLWeightTheory.derivative_digit_expansion}, and the stability of digit classes: if all
monomials $x^m$ of $F$ have the same residues $m_j\bmod P_j$, then $\partial F/\partial x_k$ has the same residues except at
$k$, where the residue $r_k$ becomes $r_k-1$ if $r_k\ge1$ and $P_k-1$ otherwise, since $(m_k-1)\bmod P_k$ is computed in this
way.
\end{itemize}
\emph{(T6).} The seven statements are
\lref{Formalism.MainNotions.LemmaBAlt.automorphicSideConclusion_of_pathBasis} (from the path basis of (T7) over
rings), \lref{Formalism.MainNotions.G.ColumnFamilies.splitting_third_part_holds} (from valuation bounds for words and
the valuation lemmas (T2)), \lref{Formalism.MainNotions.LemmaBAlt.prvB1_of_D7},
\lref{Formalism.MainNotions.LemmaBAlt.lemma_B2},
\lref{Formalism.MainNotions.LemmaBAlt.nested_sets_lemma_of_admitted} (from the joint-dimension statement and the
independence of paired paths), \lref{Formalism.MainNotions.LemmaBAlt.coefficientDependence} and
\lref{Formalism.MainNotions.LemmaBAlt.canonical_translate_jet_identity}.

\emph{(T7)} is \lref{Formalism.MainNotions.G.ColumnFamilies.admissible_polynomial_basis}.
\end{proof}

\begin{remark}\label{pf:sh-curv:tools-rem}
(1) The components (a), (c), (d) of \cref{pf:sh-curv:admitted} are part of the record $\mathrm{adm}(D)$ but are not used in
the proof of the tool statements (the component (b) is used through its own proof). (2) In (T1)(a) the denominator bound
is proved without the product formula, although it would also follow from it: by \cref{pf:sh-curv:invformula} and
Kummer's theorem each factor $p_{ab}(b-a,T)$ needs at most the denominator $p^{d_p(n,\ell(M))}$ on $M\cap V_0$
(\lref{Formalism.MainNotions.G.GLWeightTheory.invariant_projection_denominator} is a formal version of this argument,
which is not used).
\end{remark}

\subsubsection{Threshold functions for the first theorem}

\begin{definition}[The hypotheses $(\mathrm L8)$ and the first theorem]\label{pf:sh-curv:L8def}
(a) $(\mathrm{L}8)$ denotes the conjunction of the following local hypotheses: for every datum,
the \emph{second part} (the record $\mathrm{AdmB}$ of path statements described in the proof of \cref{pf:sh-curv:firstthm}
implies the closedness of the first theorem), the integrality from the metric for large constants, the extended closed form
of the column systems and two window clauses; and, independently of the datum, the Weyl volume statement, the Witt
identification, the flatness statements $\sfname{FlatTop}$ and $\sfname{FlatHH}$, the dual height basis and inclusion,
$\sfname{BePerfect}$, $\sfname{HLSpans}$, $\sfname{TestSide}$, $\sfname{LPOneEvery}$, the Cartan domain and surjectivity
statements, four clauses of the counting input (among them the signature matching and Gauss's lemma in the form of
\cref{pf:sh-curv:admitted}(d)), two threshold clauses and an instance clause. In \cref{sec:pf:linkG} these are supplied by
the local flatness results (\lref{Formalism.L2.G.local_flatness_results} and \lref{Formalism.MainNotions.G.loc8All_of}).

(b) For constants $C_0,C_H,C$ the \emph{first theorem} $\mathsf T_I(C_0,C_H,C)$ is the statement: for all
$p,n',m\in\N$, $TT\subseteq\Z/n'$, $G',N_1\in\N$, $Y,\bar Y,a,b',b_l,b_u\in\R$ and $c_{\mathrm{cl}},D_p,c_*\in\N$ satisfying
the \emph{hypotheses of type I} $\mathrm{Hyp}_I(C_0,C_H)$ (an explicit list of inequalities between these parameters,
including $1\le Y$, $(m+1)^2Y^2\le p$, the even-gap condition for $TT$ with $(G',N_1)$, $8(D_p+1)\le c_*$, lower bounds for
$D_p$ and $b'$ and the absence of odd powers of $p$ close to $a/b'$), there is a function $\kappa\ge0$ on pairs of weight
vectors $TT\to\Z$ whose mean over the boxes $I_a(m,Y,a)^{TT}\times I_b(b_l,b_u)^{TT}$ is at most $C|TT|N_1/Y$, such that
for every unramified quadratic setting (\cref{pf:linkG:hida}), all base weights $(a_w,b_w,\lambda_0)$ which are dominant
with $a_w(s)\in I_a$, $b_w(s)\in I_b$ for $s\in TT$, and all $t\ge t_I(C_0,C,\dots,a_w,b_w)$ (an explicit threshold),
$\mathsf{LP}(t,D_p,\kappa(a_w,b_w),c_*)$ holds with $r=1$ (\cref{pf:linkG:LP}).
\formal{Formalism.MainNotions.G.Loc8All}
\formal{Formalism.MainNotions.flipOneThree.Conclusion2}
\formal{Formalism.MainNotions.flipOneThree.HypI}
\formal{Formalism.MainNotions.flipOneThree.meanKap}
\formal{Formalism.MainNotions.flipOneThree.tI}
\end{definition}

\begin{theorem}[Threshold functions for the first theorem]\label{pf:sh-curv:firstthm}
Assume $(\mathrm E)$ and $(\mathrm L8)$. There are $T_0>0$ and functions $T_2,C\colon\R\to\R$ such that for every $C_0\ge
T_0$: $T_2(C_0)>0$, $C(C_0)>0$, and $\mathsf T_I(C_0,C_H,C(C_0))$ holds for every $C_H\ge T_2(C_0)$ (and every datum $D$,
which does not occur in $\mathsf T_I$).
\formallabel{Formalism.MainNotions.G.first_theorem_threshold_functions}
\formallabel{Formalism.MainNotions.G.first_theorem_uniform_thresholds}
\formallabel{Formalism.MainNotions.G.first_theorem_with_constants}
\formallabel{Formalism.MainNotions.G.first_theorem_admitted_closed}
\formallabel{Formalism.MainNotions.G.first_theorem_tools_admitted}
\formallabel{Formalism.MainNotions.G.thresholds_and}
\end{theorem}

The proof uses the following three statements on the polynomial models of the Weyl modules of $\GL_n$, proved in the
development (for $n\ge1$ and a weight $\mu$; brackets are the products of minors attached to the tableaux of shape $\mu$):
\begin{itemize}
\item (Global sections) For $R$ a field, a discrete valuation ring or $\Z$, the module of global sections $H^0_R(\mu)$ (functions
on $n\times n$ matrices over $R$ transforming under the upper unitriangular group by the character $\mu$) is spanned by the
brackets if $\mu$ is dominant, and is $0$ otherwise (\lref{Formalism.MainNotions.G.ColumnFamilies.global_sections_schur});
\item (Brackets) for dominant $\mu$, the $\Z$-span of the brackets is the set of elements of their $\Q$-span which lie in the
maximal lattice $\nabla$ attached to the canonical bracket
(\lref{Formalism.MainNotions.G.ColumnFamilies.bracket_polynomial_span});
\item (Duality) for an irreducible representation with highest weight vector $v$ and a dual highest weight vector $v'$
which is integral and primitive on $\Delta(v)$, the minimal lattice $\Delta(v)$ is the set of vectors on which every element
of the maximal lattice $\nabla(v')$ of the dual takes integral values
(\lref{Formalism.MainNotions.G.ColumnFamilies.admissible_lattice_biduality}).
\end{itemize}
These are the classical descriptions of induced modules, Weyl modules and their duality over $\Z$ and over discrete
valuation rings \cite[II.2--II.4]{Jantzen2003}.

\begin{proof}
\emph{Reduction to one datum.} If there is no datum, the statement holds vacuously with $T_0=1$ and constant functions $1$.
Otherwise fix a datum $D_0$. Since $\mathsf T_I$ does not involve the datum, it suffices to show the \emph{closedness} of
$\mathsf T_I$: there is $T_0>0$ such that for every $C_0\ge T_0$ there are $T_2,C>0$ with $\mathsf T_I(C_0,C_H,C)$ for all
$C_H\ge T_2$; choosing $T_2$ and $C$ for each $C_0$ gives the functions.

\emph{The second part.} By the second part in $(\mathrm L8)$ applied to $D_0$, closedness of $\mathsf T_I$ follows from the
record $\mathrm{AdmBE}(D_0)$. This record asserts: for every $C_c\ge0$ there is $T_0>0$ such that for every $C_0\ge T_0$
there is $T_1>0$ such that for all $C_{E1},C_{E2},C_m\ge T_1$ there is $T_2>0$ such that for all $C_H\ge T_2$ there are
constants $C_A,C_\kappa,C_t,C_s>0$ satisfying the conjunction $\mathrm{AdmB}$ of five statements: the path-pair
correspondence, two statements $B_2,B_3$ on the shares of bad and of $\kappa$-weighted pairs, the column family lower set
statement, and the record $\mathrm{AdmT}(D_0)$.

\emph{Verification of $\mathrm{AdmBE}(D_0)$.}
(i) $\mathrm{AdmT}(D_0)$ (\lref{Formalism.MainNotions.G.first_theorem_tools_admitted}) is a conjunction of twelve
groups of statements, each of which is a component of the tool statements (\cref{pf:sh-curv:tools}; namely the invariant
denominators, the integral invariants and the degree bounds, transported by the lemmas
\lref{Formalism.MainNotions.G.Joints.invariant_denominator_gives_exponent_data}, \lref{Formalism.MainNotions.G.Joints.integral_invariants_give_exponent_data}, and \lref{Formalism.MainNotions.G.Joints.degree_bounds_give_exponent_data}),
$(\mathrm E)$ itself, a component of $(\mathrm L8)$, or one of the statements (Global sections), (Brackets), (Duality) and the
graded lattice statement \lref{Formalism.MainNotions.G.ColumnFamilies.adm_eleven}.
(ii) The thresholds: the extended windows theorem gives, for $C_c\ge0$, thresholds beyond which the extended conclusions of
the column systems hold (it is the extended closed form in $(\mathrm L8)$ applied to the inequality admissions, which are
again components of the tool statements, of the window clauses and of the flatness statements of $(\mathrm L8)$;
\lref{Formalism.MainNotions.G.inequality_admissions_hold}); the integrality from the metric for large constants and
the two threshold clauses of $(\mathrm L8)$ give thresholds of the same shape. Two statements of this nested ``for all
sufficiently large constants'' form can be combined by taking the maximum of the thresholds at each stage; we take
$C_E=C_{E1}=C_{E2}$ and all five constants of the inner stage equal to $C_H$. This gives, for $C_H$ beyond the combined
threshold, constants $C_\kappa,C_t,C_s>0$ and the extended conclusions, and a constant $C_A>0$
from the first threshold clause applied to $(C_\kappa,C_s+C_\kappa)$.
(iii) The five statements of $\mathrm{AdmB}$ then follow from these conclusions by the transport lemmas
\lref{Formalism.MainNotions.G.Joints.bad_share_bound_of_extension}, \lref{Formalism.MainNotions.G.Joints.bad_share_of_extension_proposition}, \lref{Formalism.MainNotions.G.Joints.kappa_share_bound_of_extension}, and \lref{Formalism.MainNotions.G.Joints.alternative_extension_gives_bound},
and (i); for the first of them the hypotheses of type I are checked by
\lref{Formalism.MainNotions.G.Joints.weight_condition_of_integral_form} (from the $\nabla$--$\Delta$ comparison of
(T1)), \lref{Formalism.MainNotions.G.Joints.extension_gives_ratio_bound},
\lref{Formalism.MainNotions.G.Joints.fourth_extension_gives_ratio_bound} and
\lref{Formalism.MainNotions.G.Joints.residue_bound_of_extension}. The constant $C_s$ of $\mathrm{AdmB}$ is
$C_s+C_\kappa$.
\end{proof}

% ===========================================================================
\subsection{Divisibility: Gauss jets and exponential gauges}\label{sec:pf:sh-curv:div}

The two results of this subsection are used in \cref{sec:pf:reckernel} (lower bounds for the local systems at the second
real place, \cref{pf:reckernel:rowTwU}), in \cref{sec:pf:top-shape} (the orbit letter bounds in the size estimates of the
families) and in \cref{sec:pf:onto3} (chart independence of harmonic values, \cref{pf:onto3:chartindep}).

\subsubsection{Orbit letter bounds}

\begin{definition}[Local jet data and their valuation bounds]\label{pf:sh-curv:jetloc}
Let $\ell\ge5$ be a prime and $\C_\ell$ the completion of an algebraic closure of $\Q_\ell$. A \emph{local jet datum of the
first kind} $L$ consists of an even $d\ge2$, integers $N_0<N_1$, a power series $\Theta\in\C_\ell[[X_{s,a,b}]]$ in the
variables indexed by $s\in\Z/d$, $a<N_0$, $b<N_1$ (the \emph{theta series}), Young diagrams $\nu_s$ ($s\in\Z/d$) with at most
$N_0$ rows, and a set $TT\subseteq\Z/d$. Write $A_s$ and $B_s$ for the lengths of the first two rows of $\nu_s$, $|\nu_s|$ for
its size, $\gamma_\ell=\ell/(\ell^2-1)$ and $\rho_\ell=1/(\ell^2-1)$.
\begin{enumerate}[label=(\alph*)]
\item (Valuation bounds) For $B\in\R$, $\mathrm{Kval}_{\ge}(\Theta,\nu;B)$ means: for every choice $I$ of a column subset
$I_{s,c}\subseteq\{0,\dots,N_1-1\}$ of size the length of the $c$-th column of $\nu_s$, every coefficient of the polynomial
obtained from the product of minors $G_I=\prod_s\prod_c\det\bigl(\sum_aY_{s,i,a}Z_{s,a,I_{s,c}(k)}\bigr)_{i,k}$ by replacing
each monomial $Z^\beta$ by $\beta!$ times the coefficient of $X^\beta$ in $\Theta$ (where $Z_{s,a,b}$ corresponds to
$X_{s,a,b}$) has absolute value at most $\ell^{-B}$.
\item (Admissibility) $L$ is \emph{admissible} if there are $\epsilon\in\C_\ell$ with $|\epsilon|=1$ and a power series $\beta$
in variables indexed by $\Z/d$ such that: (B0) the series obtained from $\Theta$ by the substitution
$X_v\mapsto\epsilon\,\mathrm{jlog}(X_v)$, with the jet logarithm $\mathrm{jlog}(X_{s,a,b})=\sum_{e\ge0}\ell^{-e}
X_{s-2e,a,b}^{\ell^{2e}}$, factors as $U\cdot\mathrm{jlog}^*\exp(\pm\sum_{a,b}\beta(X_{\cdot,a,b}))$ with $U$ integral;
(B1) the coefficients of $\beta$ lie in the closure $W$ of the extension of $\Q_\ell$ generated by the roots of unity of
order prime to $\ell$; (B2) $\beta$ is invariant under the twists of its variables by the powers $u^{\pm\ell^s}$ of the
$(\ell^d-1)$-th roots of unity $u$; (B3) every monomial of $\beta$ has non-zero degree in the even and in the odd variables
and bidegree $\ne(1,1)$; (B4) the exponentials of the two coboundaries of $\beta$, after the jet logarithm, have integral
coefficients in $W$ and constant term $1$.
\item (Hida condition) For $a\in\R$: $\ell\le a$, $|\nu_s|\le a$ for all $s$, and $2\log_\ell a+2\le N_0-1$.
\item (Bounds) Put $D_\ell(s)=|\nu_s|-A_s$ for $s\in TT$ and $D_\ell(s)=|\nu_s|$ otherwise, and
$\mathrm{lp}_\ell(D)=\gamma_\ell\sum_sD(s)-\max\{\sum_{s,i}x_{s,i}\}$, the maximum over $x_{s,i}\ge0$ supported on odd $i<d$
with $\sum_{i<d}\ell^ix_{s,i}+\sum_{i<d}x_{s-i,i}\le D(s)$ for all $s$. With the explicit \emph{top terms} $\mathrm{top}(s)$,
\emph{base terms} $\mathrm{base}(s)$ (series over odd $e$ of positive parts of $A_s/\ell^e$, resp.\ $(A_s-B_s)/\ell^e$, minus
the capacities of $s+e$) and \emph{capacities} $W_c(s)$, the statements are:
\begin{itemize}
\item $\mathrm{ThmA}(C)$: for all $a$ with the Hida condition,
$\mathrm{Kval}_\ge\bigl(\Theta,\nu;\sum_{s\in TT}(\mathrm{top}(s)-\gamma_\ell W_c(s))+\mathrm{lp}_\ell(D_\ell)
-C\,dN_0N_1(\log_\ell a+1)\bigr)$;
\item $\mathrm{Cor}(C_6,C_d)$: for all $a$ with the Hida condition, $k\ge3$ and $w,E\colon\Z/d\to\R$ with $w>0$,
$0\le E\le4w$, $D_\ell(s)=(k-1)w(s)+E(s)$ for $s\notin TT$ and $D_\ell(s)=(k-3)w(s)+E(s)$ for $s\in TT$:
$\mathrm{Kval}_\ge\bigl(\Theta,\nu;\sum_{s\in TT}\mathrm{base}(s)+\mathrm{lp}_\ell((k-1)w)-C_6\sum_sw(s)/\ell
-\gamma_\ell\sum_{s\in TT}W_c(s)-C_d\,dN_0N_1(\log_\ell a+1)\bigr)$.
\end{itemize}
\end{enumerate}
A \emph{local jet datum of the second kind} is the analogous datum with $d=1$, one Young diagram $\nu$ and the analogous
admissibility (with a series $\beta$ in two variables); $\mathrm{ThmOne}(C)$ is the statement: for all $a\ge\ell$ with
$|\nu|\le a$, $\mathrm{Kval}_\ge\bigl(\Theta,\nu;\rho_\ell|\nu|-C\,N_0N_1(\log_\ell a+1)\bigr)$.
\formal{Formalism.BasicNotions.JetLoc}
\formal{Formalism.MainNotions.product_zeroth_first_parts.Loc}
\formal{Formalism.MainNotions.product_zeroth_first_parts.KvalGe}
\formal{Formalism.MainNotions.product_zeroth_first_parts.Loc.AdmB}
\formal{Formalism.BasicNotions.lpF}
\formal{Formalism.MainNotions.product_zeroth_first_parts.Loc.ThmA}
\formal{Formalism.MainNotions.product_zeroth_first_parts.Loc.Cor2}
\formal{Formalism.MainNotions.product_zeroth_first_parts.Loc1.ThmOne}
\end{definition}

\begin{theorem}[Orbit letter bounds with constants $(11,1)$]\label{pf:sh-curv:split}
Every admissible local jet datum of the first kind satisfies $\mathrm{ThmA}(1)$ and $\mathrm{Cor}(11,1)$, and every
admissible local jet datum of the second kind satisfies $\mathrm{ThmOne}(1)$.
\formal{Formalism.MainNotions.product_zeroth_first_parts.SplitConstantsUniformBound}
\formallabel{Formalism.MainNotions.DIV.Bridge.DworkPowerSeriesWeights.split_bound_eleven_one}
\formallabel{Formalism.MainNotions.DIV.Bridge.DworkPowerSeriesWeights.thmA_of}
\formallabel{Formalism.MainNotions.DIV.Bridge.DworkPowerSeriesWeights.thmOne_of}
\formallabel{Formalism.MainNotions.DIV.Bridge.DworkPowerSeriesWeights.corollary_split_eleven_one}
\formallabel{Formalism.MainNotions.DIV.Bridge.DworkPowerSeriesWeights.kval_mono}
\end{theorem}

The proof rests on two theorems on \emph{Gauss jets} proved in the development. For an unramified subfield $K_0\subset\C_\ell$,
a primitive $(\ell^d-1)$-th root of unity $\zeta$ and a series $\beta$, the \emph{Dwork hypotheses}
$\mathrm{BHyp}(K_0,\zeta,\beta)$ express that $\beta$ is a Dwork-type exponential series over $K_0$ with Frobenius
structure given by $\zeta$; for a substitution $g$ inverse to the entire logarithm, an integral series $U$ and a sign
$\eps$, $\theta_Z(g,U,\eps,\beta)$ is the corresponding twisted theta series, and the \emph{frame jets} of a series are the
polynomials obtained by evaluating the products of minors of (a) on it.
\begin{itemize}
\item (Theorem A for Gauss jets) Under the Dwork hypotheses, for $\ell\ge3$ and every integer $a\ge1$ with $|\nu_s|\le a$
and $2\log_\ell a+2\le N_0-1$, the coefficients of the frame jets of $\theta_Z(g,U,\eps,\beta)$ have absolute value at most
$|\ell|^{B_A(a)}$ with $B_A(a)=\sum_{s\in TT}(\mathrm{top}(s)-\gamma_\ell W_c(s))+\mathrm{lp}_\ell(D_\ell)-dN_0N_1(\log_\ell
a+1)$ (\lref{Formalism.DworkPowerSeriesWeights.thmAGauss}, proved from a transfer to a residually generic field,
the evaluation of frame jets, the existence of labels and the computation of powers of Gauss sums;
\lref{Formalism.DworkPowerSeriesWeights.thmAGauss_of}).
\item (One-letter Gauss jets) The analogous bound with exponent $\rho_\ell\sum_cp_c-N_0N_1(\log_\ell a+1)$ for $d=2$
(\lref{Formalism.DworkPowerSeriesWeights.oneLetterGauss}).
\end{itemize}
They are $\ell$-adic estimates for Gauss sums in the spirit of the Gross--Koblitz formula and Stickelberger's theorem
\cite[Ch.~6]{Washington1997}. Three bridge statements connect them with the local jet data: (Tie) the hypotheses (B1)--(B4)
imply the Dwork hypotheses for some unramified $K_0$ and primitive $\zeta$
(\lref{Formalism.MainNotions.DIV.Bridge.OneLetterJetKernel.tieB}); (Kernel) (B0) expresses the twisted theta series as
$\theta_Z(g,U,\eps,\beta)$ for suitable $g$, integral $U$ and $\eps=\pm1$
(\lref{Formalism.MainNotions.DIV.Bridge.OneLetterJetKernel.kernel_from_twist_holds}, and
\lref{Formalism.MainNotions.DIV.Bridge.OneLetterJetKernel.kernel_twist_one_letter_holds} for the second kind); (Frame)
a bound $|\ell|^{B}$ for the coefficients of the frame jets of the twisted theta series implies
$\mathrm{Kval}_\ge(\Theta,\nu;B)$ (\lref{Formalism.MainNotions.DIV.Bridge.OneLetterJetKernel.tieFrame},
\lref{Formalism.MainNotions.DIV.Bridge.OneLetterJetKernel.tieFrameOne}).

\begin{proof}
\emph{$\mathrm{ThmA}(1)$.} Let $L$ be admissible, with $\epsilon,\beta$ as in (b), and let $a$ satisfy the Hida condition.
By (Tie) we get $K_0,\zeta$ with the Dwork hypotheses, and by (Kernel) $g,U,\eps$ with $\Theta^\epsilon=\theta_Z(g,U,\eps,\beta)$,
where $\Theta^\epsilon$ is the substituted series of (B0). Put $a'=\lfloor a\rfloor\ge1$ (as $a\ge\ell>1$). The sizes are
integers $\le a$, hence $\le a'$; the column lengths of the $\nu_s$ are between $1$ and $N_0$; and $2\log_\ell a'+2\le
2\log_\ell a+2\le N_0-1$. Theorem A for Gauss jets gives the bound $|\ell|^{B_A(a')}$ for the coefficients of the frame jets
of $\Theta^\epsilon$, where the shape parameters of Theorem A are the first row lengths, the numbers of columns of length
$\ge2$ (which is $B_s$) and the sizes of the $\nu_s$; and the linear program value $F_{\mathrm{num}}$ used there equals
$\mathrm{lp}_\ell$ (\lref{Formalism.MainNotions.DIV.Bridge.DworkPowerSeriesWeights.lpF_eq_Fnum}, an identification of two
descriptions of the same feasible set). Since $|\ell|=\ell^{-1}$, $\ell>1$ and $\log_\ell a'\le\log_\ell a$, we have
$|\ell|^{B_A(a')}\le\ell^{-B_A(a)}$, and (Frame) gives $\mathrm{ThmA}(1)$.

\emph{$\mathrm{Cor}(11,1)$.} Let $a,k,w,E$ be as in $\mathrm{Cor}$. Two inequalities are used:
(i) $\mathrm{base}(s)\le\mathrm{top}(s)$ for every $s$
(\lref{Formalism.MainNotions.DIV.Bridge.DworkPowerSeriesWeights.base_term_le_top_term}): termwise, since
$(A_s-B_s)/\ell^e\le A_s/\ell^e$ and the capacities are non-negative, and both series are dominated by the geometric series
$A_s\sum_e\ell^{-e}$;
(ii) $\mathrm{lp}_\ell((k-1)w)-11\sum_sw(s)/\ell\le\mathrm{lp}_\ell(D_\ell)$ under the hypotheses on $D_\ell$
(\lref{Formalism.MainNotions.DIV.Bridge.DworkPowerSeriesWeights.local_bound_eleven_comparison}, a comparison of the
two linear programs obtained from the bound \lref{Formalism.DworkPowerSeriesWeights.OddLevelGainBounds.normCor7} on the
gain of the odd levels, with $\epsilon=\frac12$).
Together they show that the bound of $\mathrm{Cor}(11,1)$ is at most the bound of $\mathrm{ThmA}(1)$. Since
$\mathrm{Kval}_\ge(\Theta,\nu;B)$ implies $\mathrm{Kval}_\ge(\Theta,\nu;B')$ for $B'\le B$ (as $\ell^{-B}\le\ell^{-B'}$),
$\mathrm{ThmA}(1)$ gives $\mathrm{Cor}(11,1)$.

\emph{$\mathrm{ThmOne}(1)$.} The same argument with the one-letter versions: (Tie) for $d=2$, (Kernel) for the second kind,
the one-letter Gauss jets with $a'=\lfloor a\rfloor$ (the column lengths sum to $|\nu|\le a'$), and (Frame) for the second
kind; $\log_\ell a'\le\log_\ell a$ again gives the bound with $a$.
\end{proof}

\begin{remark}\label{pf:sh-curv:split-rem}
The constants $(11,1)$ are explicit: the factor $1$ in front of $dN_0N_1(\log_\ell a+1)$ comes directly from Theorem A for
Gauss jets, and the constant $11$ from the comparison (ii) of the linear programs.
\end{remark}

\subsubsection{Exponential gauges}

\begin{definition}[Exponential charts and gauges]\label{pf:sh-curv:expdef}
Let $E$ be a CM field with a host structure (whose involution is the complex conjugation of $E$), $\mathcal X$ a moduli space
of the development attached to a datum and a level, with universal group scheme $G\to\mathcal X$, its action of $\OO_E$, and
the rigidified automorphic line bundle $\mathcal L$ on $G$; let $c$ be a coefficient datum and $S$ a $c.F$-algebra containing
$\Q$. Let $y$ be a \emph{framed point} of $\mathcal X$ over $S$ (a point with a frame $u$ of the tangent space indexed by the
set $\Lambda_1$ of $d_L$ letters) and $J_m=S[x_1,\dots,x_{d_L}]/(x)^{m+1}$.
\begin{enumerate}[label=(\alph*)]
\item A \emph{chart} $C$ at $y$ consists of a formal group law $F$ of dimension $d$ over $S$, compatible maps $x^C_m$ from
the spectrum of the jets of order $m$ in $d$ variables (over $y$) to $G$, the action of $\OO_E$ on the coordinates,
generators $e^C_m$ of $\mathcal L$ along $x^C_m$ and further structure (a principal element $w_P$ and frames for the adjoint
relations); it is \emph{truncation compatible} if $e^C_{m+1}$ restricts to $e^C_m$ for all $m$.
\item An \emph{exponential} of $C$ is a $d$-tuple $\varphi$ of power series in $d_L$ variables without constant term, whose
linear terms are given by the frame, and which is a homomorphism from the additive formal group to $F$:
$\varphi(s+t)=F(\varphi(s),\varphi(t))$. It defines maps $x_m=x^C_m\circ\varphi\colon\mathrm{Spec}\,J_m\to G$.
\item The family $(x_m)$ is an \emph{exponential chart} if $x_m(0)=1$, the $x_m$ are compatible with truncation,
$x_m(s+t)=x_m(s)x_m(t)$, the action of $a\in\OO_E$ on $G$ corresponds to the scaling of the coordinates by the images of $a$
under the embeddings attached to the letters, and the tangent vectors of $x_1$ are given by the frame $u$.
\item An \emph{exponential gauge} for an exponential chart $(x_m)$ is a family $e=(e_m)$ of generators of $x_m^*\mathcal L$,
compatible with truncation, equal at $m=0$ to the rigidification, and satisfying the \emph{cocycle condition}: for every
$a\in\OO_E$ there is a generator $\beta_a$ of $(1+a\bar a)^*\mathcal L\otimes a^*\mathcal L^{-2}\otimes\mathcal L^{-1}
\otimes(a\bar a)^*\mathcal L^{-1}$ on $G$, whose restriction to the unit section is the $(-3)$-rd power of the
rigidification, such that $e_m\circ(1+a\bar a)=x_m^*\beta_a\cdot(e_m\circ a)^2\cdot e_m\cdot(e_m\circ a\bar a)$ for all $m$,
where $e_m\circ b$ denotes the pull-back of $e_m$ along the action of $b$ on the coordinates.
\item The gauge $e$ is \emph{the gauge of $(C,\varphi)$} if $e_m=u_m\cdot\varphi^*e^L_m$ for all $m$, where $e^L_m$ is the
generator of $\mathcal L$ along $x^C_m$ built from the chart and $u_m\in J_m^\times$ is the truncation of
$\exp(-\varphi^*h_L)$, $h_L$ the logarithmic series of the chart at the base point.
\end{enumerate}
The \emph{exponential action} statement $(\mathrm{EA})$ asserts: for every chart $C$ at a framed point and every exponential
$\varphi$ of $C$, the action of $a\in\OO_E$ on the chart coordinates, composed with $\varphi$, is the linear scaling of the
variables of $\varphi$ by the images of $a$ under the embeddings attached to the letters.
\formal{Formalism.MainNotions.ModuliIntegralCharts.H1.IsExpOf}
\formal{Formalism.L2.Pj.IsExpChart}
\formal{Formalism.L2.Pj.IsExpGauge}
\formal{Formalism.MainNotions.ModuliIntegralCharts.H1.IsExpGaugeOf}
\formal{Formalism.MainNotions.ModuliIntegralCharts.H1.FrTruncC}
\formal{Formalism.MainNotions.ModuliIntegralCharts.H1.ExpActOfSt}
\end{definition}

\begin{theorem}[Exponential gauges]\label{pf:sh-curv:expgauge}
Assume $(\mathrm{EA})$. Let $C$ be a truncation compatible chart at a framed point $y$ over $S$ and $\varphi$ an exponential
of $C$ such that $(x^C_m\circ\varphi)_m$ is an exponential chart. Then there is an exponential gauge for
$(x^C_m\circ\varphi)_m$ which is the gauge of $(C,\varphi)$.
\formal{Formalism.MainNotions.ModuliIntegralCharts.H1.ExpGaugeOfSt}
\formallabel{Formalism.MainNotions.ModuliIntegralCharts.H1.ExpGauge.expGaugeOf_of}
\end{theorem}

\begin{proof}
\emph{Construction.} By \lref{Formalism.MainNotions.ModuliIntegralCharts.H1.ExpGA.exists_eG} there are generators
$\tau_m$ of $\mathcal L$ along $x^C_m$ corresponding to $e^L_m$ and units $u_m\in J_m^\times$ equal to the truncations of
$\exp(-\varphi^*h_L)$ (this uses that $\varphi$ has no constant term), such that $e_m:=u_m\cdot\varphi^*\tau_m$ are
generators and $e$ is the gauge of $(C,\varphi)$; here $\varphi^*\tau_m$ is the pull-back along the map of jets induced by
$\varphi$.

\emph{Truncation and rigidity.} The $e_m$ are compatible with truncation, because the $e^C_m$ are (truncation compatibility of
$C$) and the units $u_m$ are truncations of one power series
(\lref{Formalism.MainNotions.ModuliIntegralCharts.H1.ExpGB.trunc_eG}); and $e_0$ is the rigidification
(\lref{Formalism.MainNotions.ModuliIntegralCharts.H1.ExpGB.rig_eG}, using again $\varphi(0)=0$).

\emph{The cocycle condition.} Let $a\in\OO_E$. Since the involution of $\OO_E$ is the host conjugation, $\bar{\bar a}=a$.
The frame of the chart provides generators $w_a,w_{\bar a}$ satisfying the adjoint relations for $a$ and $\bar a$
(\lref{Formalism.MainNotions.ModuliIntegralCharts.H1.ExpGD.exists_adRel}). From them and the principal element $w_P$ of
the chart, \lref{Formalism.MainNotions.ModuliIntegralCharts.H1.ExpGE.exists_beta} constructs a generator $\beta_a$ of the
line bundle of (d), with restriction to the unit section equal to the $(-3)$-rd power of the rigidification, and with
explicit values in terms of a quotient $B_q(w_P,a,w_a,w_{\bar a})$. Finally the identity of (d) for $e$ and $\beta_a$ holds
for every $m$ (\lref{Formalism.MainNotions.ModuliIntegralCharts.H1.ExpGL.coc_chart}): it is checked on values, using
$(\mathrm{EA})$ (which turns the action of $a$, $a\bar a$ and $1+a\bar a$ on the chart into scalings of the variables of
$\varphi$), the values of the phase factors of $h_L$ at one and two points
(\lref{Formalism.MainNotions.ModuliIntegralCharts.H1.ExpGJPh1.phOneVal_holds} and \lref{Formalism.MainNotions.ModuliIntegralCharts.H1.ExpGJ.phTwoVal_holds}),
the values of the adjoint generators (\lref{Formalism.MainNotions.ModuliIntegralCharts.H1.ExpGJ.adVal}), a unit identity
(\lref{Formalism.MainNotions.ModuliIntegralCharts.H1.ExpGU.unitId}) and the value of $B_q$ at the origin
(\lref{Formalism.MainNotions.ModuliIntegralCharts.H1.ExpGO.bqOrigin}).
\end{proof}

% ===========================================================================
\subsection{Archimedean inputs at built places}\label{sec:pf:sh-curv:arch}

The six results of this subsection are inputs of the carrier statement of Theorem~S (\cref{sec:pf:thmS-carrier}) and of
the leaf and pinning statements (\cref{sec:pf:thmS-leaves,sec:pf:thmS}). We use the archimedean objects of
\cref{sec:pf:thmS-carrier:objects}, which we recall.

\begin{definition}[Archimedean objects]\label{pf:sh-curv:archdef}
For integers $k\ge1$, $r\ge0$ let $\mathrm{Mat}_{k,r}=M_{2r+k+1}(\C)$ and
$G_{k,r}=\Uni(H_{k,r})$; a real place datum $B$ (\cref{pf:thmS-carrier:realplace}) consists of vectors $z,z'$ in the
anisotropic part with $T(z)=-\Re H_{k,r}(z,z)>0$ and $T(z')>0$, unitary representations $\pi,\pi'$ on $W_p,W_q$ of the
compact groups $U_W(z),U_W(z')$ (the elements of $G_{k,r}$ fixing $z$, resp.\ $z'$, and the isotropic basis vectors
$e_a,e^*_a$, $a<r$), Haar measures on $G_{k,r}$, $U_W(z)$, $U_W(z')$ and on the unipotent group $N=N_{k,r}$, a section
and a cut-off $\chi$. We write $g_N(g)=(\sum_{i,j}|g_{ij}|^2)^{1/2}$; it is submultiplicative and $g_N\ge1$ on $G_{k,r}$.
For $p=(h,n,h',n')\in\mathcal S=U_W(z)\times N\times U_W(z')\times N$, with the product measure $\mu_{\mathcal S}$, put
\[
s_p(g)=h^{-1}n^{-1}g\,n'h',\qquad(\rho(p)F)(g)=\pi(h)\,F(s_p(g))\,\pi'(h')^{-1},\qquad
\xi(p)=e^{i(\lambda_z(n)-\lambda_{z'}(n'))},
\]
where $\lambda_z,\lambda_{z'}$ are the real characters of $N$ attached to $z,z'$; the \emph{Bessel transform} is
$J_F(g)=\int_{\mathcal S}\xi(p)(\rho(p)F)(g)\,d\mu_{\mathcal S}(p)$. A function $F\colon\mathrm{Mat}_{k,r}\to\Hom(W_q,W_p)$
is of \emph{class $\mathrm{CC}(A,m)$} if it is $C^m$ near $G_{k,r}$ and $g_N(g)^A\|(D_uD'_vF)(g)\|$ is bounded on $G_{k,r}$ for
all words $u,v$ of total length $\le m$ in unit elements of the Lie algebra, $D_u,D'_v$ the left and right derivatives.
\formal{Formalism.FourierDecayAlt.B.RealPlaceData}
\formal{Formalism.FourierDecayAlt.B.rho}
\formal{Formalism.FourierDecayAlt.B.xiTau}
\formal{Formalism.FourierDecayAlt.B.JF}
\formal{Formalism.FourierDecayAlt.B.InCC}
\end{definition}

\subsubsection{Convergence of orbital integrals}

\begin{definition}[Relevant monomial orbits and normalised cuts]\label{pf:sh-curv:orbdef}
Let $B$ be a real place datum and $\gamma\in G_{k,r}$. Let $\widehat Z(\gamma)$ be the $(r+1)\times(r+1)$ matrix of the
hermitian pairings $H_{k,r}(\gamma f'_b,f_a)$, where $f_a=f'_a=e_a$ for $a<r$, $f_r=z$ and $f'_r=z'$, and
$G_f(\gamma)=|\det\widehat Z(\gamma)|^2-T(z)T(z')|\det\widehat Z_r(\gamma)|^2$, $\widehat Z_r$ the leading $r\times r$ minor.
Let $R_\gamma=\{\varrho\in N:\gamma^{-1}\varrho\gamma\in N\}$; $\gamma$ is \emph{relevant} if
$\psi_z(\varrho)=\psi_{z'}(\gamma^{-1}\varrho\gamma)$ for all $\varrho\in R_\gamma$, where $\psi_z,\psi_{z'}$ are the
characters of $N$ attached to $z,z'$. $\widehat Z(\gamma)$ is \emph{monomial} with permutation $w$ and entries $d$ if
$\widehat Z(\gamma)_{w(a),a}=d_a\ne0$ and all other entries vanish; then $R_\gamma$ carries the measure $\mu_R$, the image of
Lebesgue measure under an explicit parametrisation of $R_\gamma$ by the coordinates $(i,j)$, $j<i$, $w^{-1}(i)<w^{-1}(j)$. A
\emph{normalised cut} is a measurable $\theta\colon\mathcal S\to[0,\infty)$ with
$\int_{R_\gamma}\theta(\varrho\cdot p)\,d\mu_R(\varrho)=1$ for all $p\in\mathcal S$, where
$\varrho\cdot(h,n,h',n')=(h,\varrho n,h',\gamma^{-1}\varrho\gamma n')$.
\formal{Formalism.QuantumDeterminants.Zhat}
\formal{Formalism.QuantumDeterminants.Gf}
\formal{Formalism.FourierDecayAlt.C6.IsRelevant}
\formal{Formalism.FourierDecayAlt.D.muR}
\formal{Formalism.FourierDecayAlt.D.IsNormCut}
\formal{Formalism.FourierDecayAlt.D.orbIntegrand}
\end{definition}

\begin{theorem}[Convergence of orbital integrals]\label{pf:sh-curv:orb}
There is a constant $C_{\mathrm{orb}}$ with the following property. Let $B$ be a real place datum, $\gamma\in G_{k,r}$
relevant with $\widehat Z(\gamma)$ monomial and $G_f(\gamma)>0$, $A\ge C_{\mathrm{orb}}(r+1)^2(r(2r+k+1)+1)$, and $F$ of
class $\mathrm{CC}(A,1)$. Then there is a normalised cut $\theta$ such that $p\mapsto\theta(p)\,\xi(p)\,(\rho(p)F)(\gamma)$ is
$\mu_{\mathcal S}$-integrable.
\formal{Formalism.FourierDecayAlt.D.OrbJConv}
\formallabel{Formalism.FourierDecayAlt.P.orbital_integral_converges}
\formallabel{Formalism.FourierDecayAlt.P.Dx3.d_orbJConv}
\end{theorem}

\begin{proof}
\emph{Constants.} By the properness of the polynomial maps attached to monomial orbits
(\lref{Formalism.FourierDecayAlt.P.polynomial_map_proper}) there is $c_1\in\N$ such that for every $\gamma$ as in the
theorem there is $C_P$ with
\begin{equation}\label{eq:pf:sh-curv:proper}
g_N(\sigma)+g_N(s')\le C_P\,g_N(\sigma^{-1}\gamma s')^{c_1(r+1)^2}
\end{equation}
for all $\sigma$ in the section set of $\gamma$ and all $s'$ in the stabiliser set of $z'$. Put $C_{\mathrm{orb}}=16(c_1+1)$,
$n=2r+k+1$, $c=(c_1+1)(r+1)^2$, $b=4(rn+1)$ and $M=A/c-2b$. The hypothesis on $A$ gives $A/c\ge16(rn+1)=4b$; since
$(r+1)^2\le rn+1$, we get $M\ge2b>2(r+1)^2$ and $M\ge0$.

\emph{The cut.} By \lref{Formalism.FourierDecayAlt.P.Dx2.Rad.haar_good}, $\mu_R$ is a Haar measure on $R_\gamma$ (in
particular left invariant). Since $M>2(r+1)^2$, which exceeds the real dimension of the parameter space of $R_\gamma$, and the
parametrisation has polynomial growth, $x\mapsto g_N(x)^{-M}$ is $\mu_R$-integrable
(\lref{Formalism.FourierDecayAlt.P.Dx3.intR}), and $c_R=\int g_N^{-M}d\mu_R>0$
(\lref{Formalism.FourierDecayAlt.P.Dx3.om_pos}). Put $\Omega(n)=\int g_N(\varrho n)^{-M}d\mu_R(\varrho)$ and
$\theta(h,n,h',n')=g_N(n)^{-M}/\Omega(n)$. By left invariance $\Omega(\varrho n)=\Omega(n)$, so
$\int\theta(\varrho\cdot p)\,d\mu_R(\varrho)=\Omega(n)/\Omega(n)=1$: $\theta$ is a normalised cut
(\lref{Formalism.FourierDecayAlt.P.Dx3.cut_norm}).

\emph{Pointwise bounds.} From $F\in\mathrm{CC}(A,1)$ (words of length $0$) there is $c_F\ge0$ with $\|F(g)\|\le
c_Fg_N(g)^{-A}$ on $G_{k,r}$. The groups $U_W(z),U_W(z')$ are compact, so $g_N\le K_W$ on $U_W(z)$ and $g_N\le K'_W$ on
$U_W(z')$ for some $K_W,K'_W\ge1$ (\lref{Formalism.FourierDecayAlt.P.Dx3.UW_bound}). Let $p=(h,n,h',n')\in\mathcal S$
(this holds for $\mu_{\mathcal S}$-almost every point). Then:
\begin{enumerate}[label=(\alph*)]
\item $\|(\rho(p)F)(\gamma)\|\le c_Fg_N(s_p(\gamma))^{-A}$, since $\pi,\pi'$ are unitary, and $|\xi(p)|=1$;
\item $n^{-1}\gamma n'=h\,s_p(\gamma)\,h'^{-1}$, so $g_N(n^{-1}\gamma n')\le K_WK'_W\,g_N(s_p(\gamma))$;
\item (Geometry, \lref{Formalism.FourierDecayAlt.P.Dx3.geom}) by \eqref{eq:pf:sh-curv:proper} there are $\varrho\in
R_\gamma$, $\sigma\in G_{k,r}$ and $X$ with $n=\varrho\sigma$, $g_N(\sigma)\le X\le C_Pg_N(n^{-1}\gamma n')^c$ and
$g_N(n')\le g_N(\gamma^{-1})g_N(\gamma)\,g_N(n)X^2$;
\item $\Omega(n)\ge c_R\,g_N(\sigma)^{-M}$, by $\Omega(\varrho\sigma)=\Omega(\sigma)$ and
$g_N(x\sigma)^{-M}\ge g_N(x)^{-M}g_N(\sigma)^{-M}$ (\lref{Formalism.FourierDecayAlt.P.Dx3.om_lower}); hence
$\theta(p)\le g_N(n)^{-M}g_N(\sigma)^M/c_R$.
\end{enumerate}
Combining (a)--(d) with $M=A/c-2b$ and $4b\le A/c$ (an elementary manipulation of powers,
\lref{Formalism.FourierDecayAlt.P.Dx3.key_real}) gives
\[
\theta(p)\,\|\xi(p)(\rho(p)F)(\gamma)\|\le K\,g_N(n)^{-(A/c-3b)}\,g_N(n')^{-b}
\]
with $K=c_FK_W^A(K'_W)^AC_P^{A/c}\bigl(g_N(\gamma^{-1})g_N(\gamma)\bigr)^b/c_R$.

\emph{Integrability.} Both exponents are at least $b=4(rn+1)$, and for such exponents
$p\mapsto g_N(n)^{-a}g_N(n')^{-b}$ is $\mu_{\mathcal S}$-integrable
(\lref{Formalism.FourierDecayAlt.P.Dx3.maj_int}: the measures on the compact groups $U_W$ are finite, and
$\int_Ng_N^{-a}d\mu_N<\infty$ for $a\ge4(rn+1)$). The integrand is measurable
(\lref{Formalism.FourierDecayAlt.P.Dx3.integrand_aesm}, using the continuity of $F$ on $G_{k,r}$), so it is
integrable by domination.
\end{proof}

\subsubsection{Radial spaces, words and top lifts}

\begin{definition}[Radial spaces and actual-operator data]\label{pf:sh-curv:raddef}
Let $P=(k,r,T,T',u,u')$ be a parameter (\cref{sec:pf:thmS-carrier:built-defs}) and $B$ a real place datum with $T(z)=T$,
$T(z')=T'$. The \emph{radial domain} $\mathrm{Dom}_R$ is the set of pairs $(c,g)$ with $c\in(\C^\times)^r$ and $g$ in the
level-$r$ subgroup $G^{(r)}_{k,r}$ (the elements of $G_{k,r}$ fixing $e_a,e^*_a$, $a<r$) with $|\zeta_r(g)|^2>TT'$, where
$\zeta_r(g)$ is the last leading minor. The \emph{radial space} $\mathcal R=\mathcal R(P,B)$ consists of the functions
$\Phi\colon(\C^\times)^r\times\mathrm{Mat}_{k,r}\to\Hom(W_q,W_p)$ vanishing off $\mathrm{Dom}_R$, locally the restrictions of
smooth functions there, and satisfying $\Phi(c,hgh')=\pi(h)\Phi(c,g)\pi'(h')$ for $h\in U_W(z)$, $h'\in U_W(z')$. An
\emph{actual-operator datum} at level $l_0$ consists of the actual operators $X$ on $\mathcal R$ (radial versions of the
raising, lowering, Capelli and multiplication operators), a Haar measure $\mu$ on $G^{(r)}_{k,r}$, a level datum $L$ of the
operator algebra of $\mathcal R$ and a start datum $\mathrm{st}$ (quadratic operators $QQ(Y)$, a function $q(Y)$, operators
$x(Y),\xi(Y)$, and an invertibility set $\mathrm{inv}$), such that $B$ is a place for $P$, the structure operators of $L$ are
those of $X$ (an \emph{actual level}), $\mathrm{st}.QQ=\pm X.QQ$ with the sign $(-1)^{k-l_0}$, $\mathrm{st}.q$ is the
explicit polynomial $q_C$ of $P$, $\mathrm{inv}$ is the set of $Y$ at which no type polynomial $x_C(\rho,Y)$ vanishes, and $L$,
$\mathrm{st}$ satisfy the star and start relations.
\formal{Formalism.FourierDecayAlt.Q.DomR}
\formal{Formalism.FourierDecayAlt.Q.RadSp}
\formal{Formalism.FourierDecayAlt.Q.ActualOps}
\formal{Formalism.FourierDecayAlt.P5.IsData}
\end{definition}

\begin{definition}[Words, cut words and top lifts]\label{pf:sh-curv:topliftdef}
Let $P$, $B$, $(X,\mu,L,\mathrm{st})$ and $l_0$ be as in \cref{pf:sh-curv:raddef}. A \emph{letter} is a pair $(b,Y)$ with
$b\in\{0,1\}$ and $Y\in\R$; a list $w$ of letters is a \emph{cut word} at level $l_0$ if all its letters lie on the cut of
$P$ at level $l_0$ and $|Y-Y'|>2$ for distinct letters. A word $w$ defines through $(L,\mathrm{st})$ a \emph{frame}
$\mathbf F_w$ and a \emph{slope test product} $b_T(w)$, matrices of operators on $\mathcal R$ indexed by
$\{0,1\}^{|w|}$. A \emph{compactly supported coefficient} is a smooth compactly supported $F$ on $G_{k,r}$ with all growth
bounds; $\Phi\in\mathcal R$ is a \emph{radial representative} of $F$ if it agrees with the restriction of $J_F$ to
$\mathrm{Dom}_R$.
\begin{enumerate}[label=(\alph*)]
\item A \emph{top lift} is a family $\Theta(w,e,F)$ of functions on $\mathrm{Mat}_{k,r}$, for words $w$,
$e\in\{0,1\}^{|w|}$ and functions $F$, with $\Theta(\emptyset,e,F)=F$. It is a top lift \emph{of the datum} if for every
cut word $w$, compactly supported coefficient $F$ with radial representative $\Phi$, and $e$, the $e$-th component of
$(\mathbf F_w\,b_T(w))$ applied to the first state vector of $\Phi$ is a radial representative of $\Theta(w,e,F)$.
\item $\Theta$ is \emph{of class $D^+$ with constant $C_b$} if for every cut word $w$ and every $e$ there are a
differential operator $\mathcal D$ with bounded fibre coefficients, of order $o\le C_b|w|(2r+k+1)$ and degree
$\kappa\le C_b(r+1)o$ in the class $D^+$, such that $\Theta(w,e,F)=\mathcal DF$ on the domain of $\mathcal D$ for all $F$.
\end{enumerate}
\formal{Formalism.FourierDecayAlt.KB.CutWord}
\formal{Formalism.FourierDecayAlt.Top.TopLift}
\formal{Formalism.FourierDecayAlt.KB.IsTopLift2}
\formal{Formalism.FourierDecayAlt.KB.InDDplus2}
\end{definition}

\begin{theorem}[Top lifts exist]\label{pf:sh-curv:toplift}
Let $C_b=11$. For every parameter $P$, real place datum $B$ and actual-operator datum $(X,\mu,L,\mathrm{st})$ at a level
$1\le l_0\le k$ there is a top lift $\Theta$ of the datum which is of class $D^+$ with constant $C_b$.
\formal{Formalism.FourierDecayAlt.TL.TopLiftExists3}
\formallabel{Formalism.FourierDecayAlt.P.top_lift_exists}
\end{theorem}

\begin{proof}
For $o,\kappa\in\N$ say that an operator $a$ on $\mathcal R$ \emph{lifts with $(o,\kappa)$} if there is a differential
operator $\mathcal D$ of order $\le o$ and degree $\le\kappa$ in the class $D^+$, with bounded fibre coefficients, such that
$a\Phi$ is a radial representative of $\mathcal DF$ whenever $\Phi$ is a radial representative of a compactly supported
coefficient $F$.

\emph{Step 1: the lift system.} The operators that lift form a \emph{lift system}:
(i) lifting is monotone in $(o,\kappa)$, $0$ lifts, and sums, scalar multiples and products lift, with orders and degrees
adding under products (\lref{Formalism.FourierDecayAlt.P.lift_closed_under_operations}, whose main ingredient is the
composition of differential operators of class $D^+$,
\lref{Formalism.FourierDecayAlt.P.compose_differential_operators});
(ii) the identity and the multiplication operators of the level datum lift with order $0$
(\lref{Formalism.FourierDecayAlt.P.lift_of_functions});
(iii) the operators $x(Y)$ of the start datum lift (\lref{Formalism.FourierDecayAlt.P.lift_of_x}), and for $Y\in
\mathrm{inv}$ the operator $\xi(Y)$ is a polynomial in $x(Y)$ (\lref{Formalism.FourierDecayAlt.P.actual_algebra_facts}),
so it lifts by (i);
(iv) the Capelli operators of the start datum coincide with the actual Capelli operators
(\lref{Formalism.FourierDecayAlt.P.actual_algebra_facts}) and lift
(\lref{Formalism.FourierDecayAlt.P.lift_capelli_operators}).

\emph{Step 2: lifts of words.} By the recursion over words (\lref{Formalism.FourierDecayAlt.TL.wordLiftOfParts}, from
\lref{Formalism.FourierDecayAlt.P.word_recursion} and the lift of single letters
\lref{Formalism.FourierDecayAlt.P.letter_lift_of_parts}, built from the basic lift
\lref{Formalism.FourierDecayAlt.TL.basicLift} and the shape statements
\lref{Formalism.FourierDecayAlt.P.l_shape} and \lref{Formalism.FourierDecayAlt.P.m_shape_of_l_shape}), for every word $w$ which is
admitted by the start datum and every $e$, the $(e,1)$-entry of $\mathbf F_w\,b_T(w)$ lifts with
$o=|w|(2r+k+1)$ and $\kappa=11|w|(r+1)(2r+k+1)$. Cut words are admitted: their positions are distinct, lie on the cut and
are pairwise far apart, and the start datum satisfies the interlacing and positivity conditions on the cut
(\lref{Formalism.FourierDecayAlt.P.cut_admissible}).

\emph{Step 3: the top lift.} Put $\Theta(\emptyset,e,F)=F$ and, for $w\ne\emptyset$, $\Theta(w,e,F)=\mathcal D_{w,e}F$, where
$\mathcal D_{w,e}$ is the operator of Step~2 if $w$ is a cut word and the identity operator
(\lref{Formalism.FourierDecayAlt.P.identity_operator}) otherwise. Since the first state vector of $\Phi$ has only one
non-zero component, the $e$-th component of $\mathbf F_wb_T(w)$ applied to it is the $(e,1)$-entry applied to $\Phi$, so
$\Theta$ is a top lift of the datum (for $w=\emptyset$ the frame and the slope test product are the identity). For a cut word
$w\ne\emptyset$ the bounds of Step~2 give $o=|w|(2r+k+1)\le11|w|(2r+k+1)$ and $\kappa=11(r+1)o$; for $w=\emptyset$ one takes the
operator of order and degree $0$. Hence $\Theta$ is of class $D^+$ with constant $11$.
\end{proof}

\subsubsection{Chart equations}

\begin{definition}[Charts and the Bessel system]\label{pf:sh-curv:chartdef}
Let $k\ge1$ and $z,z'$ be anisotropic vectors with $T(z),T(z')>0$; let $\Omega^+\subset\mathrm{Mat}_{k,r}$ be the open
domain of $(z,z')$ (\cref{pf:thmS-carrier:transferdef}). A \emph{chart} of dimension $d$ is a pair of smooth maps
$\kappa\colon O\to U$, $\Phi\colon U\to O$ between open sets $O\subseteq\R^d$ and $U\subseteq\mathrm{Mat}_{k,r}$ with
$\kappa(O)\subseteq G_{k,r}\cap\Omega^+$, $\Phi\circ\kappa=\mathrm{id}_O$ and $\kappa(\Phi(g))=g$ for $g\in U\cap G_{k,r}$. A
function $u\colon\mathrm{Mat}_{k,r}\to\C$ satisfies the \emph{Bessel system} of $(z,z')$ if it is of class $C^{2r+k+1}$ on
$\Omega^+$ at level $0$ and, on the open domain, it is an eigenfunction of the right and left derivatives along the Lie
algebra of $N$, of the left Capelli operators for every frame of $z$ and the right Capelli operators for every frame of
$z'$, and of the Casimir-type operators $\mathrm{CV}(Y)$, $Y\in\R$.
\formal{Formalism.FourierDecayAlt.P.CutMP.SmoothTransferRegularity.GEqs}
\end{definition}

\begin{theorem}[Chart equations]\label{pf:sh-curv:charteq}
Let $k\ge1$, $z,z'$ as above, $(\kappa,\Phi)$ a chart of dimension $d$ and $u$ a solution of the Bessel system. Then there
are finitely many scalar differential operators $P_j$ on $O$, with smooth coefficients and of order $\le2r+k+1$, such that
$u\circ\kappa$ is of class $C^{2r+k+1}$ on $O$, $P_j(u\circ\kappa)=0$ on $O$ for all $j$, and for every $x\in O$ and every
$\xi\in\R^d\smallsetminus\{0\}$ at which all principal symbols $P_j(x,\xi)$ vanish there is a non-zero nilpotent $\Xi$ in the
Lie algebra of $G_{k,r}$ with $\Re\tr(\Xi X)=0$ for all $X$ in the Lie algebras of $U_W(z)$ and $N$, and
$\Re\tr(\Xi\,\kappa(x)X\kappa(x)^{-1})=0$ for all $X$ in the Lie algebras of $U_W(z')$ and $N$.
\formal{Formalism.FourierDecayAlt.P.CutMP.SmoothTransferRegularity.ChartEqs}
\formallabel{Formalism.FourierDecayAlt.P.CutMP.SmoothTransferRegularity.chart_equations_holds}
\formallabel{Formalism.FourierDecayAlt.P.CutMP.SmoothTransferRegularity.chartEqs_of_capJoints}
\formallabel{Formalism.FourierDecayAlt.P.CutMP.SmoothTransferRegularity.chartEqs_of_J}
\formallabel{Formalism.FourierDecayAlt.P.CutMP.SmoothTransferRegularity.chartSysWOfParts_of}
\end{theorem}

Thus the characteristic directions of the system in the chart correspond to nilpotent elements annihilating the
infinitesimal symmetries of $u$ on both sides; this is the form in which the system is used to prove that the Bessel
transforms are smooth on the open domain (\cref{pf:thmS-carrier:openQQ}(T1)). It is the analogue, for the present
two-sided Bessel system, of the classical fact that the characteristic variety of a system containing all Casimir operators
is contained in the nilpotent cone.

\begin{proof}
\emph{The words.} Three statements on \emph{Capelli words} (products of Capelli operators and of right and left
Capelli operators) express their action on functions near a point through the derivatives of the function and the exponential
coordinates; they are assembled
(\lref{Formalism.FourierDecayAlt.P.CutMP.SmoothTransferRegularity.capWords_all_of}) from the local representative of
derivatives, the right and left exponential evaluations
(\lref{Formalism.FourierDecayAlt.P.CutMP.SmoothTransferRegularity.right_exponential_evaluation} and \lref{Formalism.FourierDecayAlt.P.CutMP.SmoothTransferRegularity.left_exponential_evaluation}),
the expression of the Capelli vectors and operators through generators, the degree and symbol of exponentials, the
coefficient of determinant shifts
(\lref{Formalism.FourierDecayAlt.P.CutMP.SmoothTransferRegularity.determinant_shift_coefficient}) and coefficient
interpolation.

\emph{The equations.} Given these, \lref{Formalism.FourierDecayAlt.P.CutMP.SmoothTransferRegularity.chartEqs_of_J}
produces the operators $P_j$. Fix a basis $X_1,\dots,X_q$ of the Lie algebra of $N$ and frames of $z$ and $z'$
(\lref{Formalism.FourierDecayAlt.P.CutMP.SmoothTransferRegularity.frame_existence}). The family consists of the
first-order operators $V_R(X_i)-c_R(i)$ and $V_L(X_i)-c_L(i)$ (the right and left derivatives along $X_i$ written in the chart,
minus the eigenvalues of $u$), and of operators $P^{\mathrm{CV}}_j$ ($1\le j\le2r+k+1$), $P^{R}_j$ and $P^L_j$ ($1\le j\le k$) of
order $j$ obtained from the Casimir-type and the right and left Capelli equations; the latter are produced from the Capelli
words, the chart relations and the passage from vector operators to symbol operators
(\lref{Formalism.FourierDecayAlt.P.CutMP.SmoothTransferRegularity.right_chart_holds}, \lref{Formalism.FourierDecayAlt.P.CutMP.SmoothTransferRegularity.vector_to_symbol_operator}, \lref{Formalism.FourierDecayAlt.P.CutMP.SmoothTransferRegularity.symbol_operator_field_holds}, \lref{Formalism.FourierDecayAlt.P.CutMP.SmoothTransferRegularity.symbol_operator_lift_holds}, and \lref{Formalism.FourierDecayAlt.P.CutMP.SmoothTransferRegularity.symbol_operator_linear_holds}).
The chart relations also give the regularity of $u\circ\kappa$ and the vanishing $P_j(u\circ\kappa)=0$.

\emph{The characteristic set.} Let $\xi\ne0$ be characteristic at $x$. By the covariance of the chart
(\lref{Formalism.FourierDecayAlt.P.CutMP.SmoothTransferRegularity.chart_covariance_holds} and \lref{Formalism.FourierDecayAlt.P.CutMP.SmoothTransferRegularity.left_covariance_chart_holds})
there is $\Xi_0\ne0$ in the Lie algebra of $G_{k,r}$ such that the principal symbol at $(x,\xi)$ of the right derivative along
$X$ is $\Re\tr(\Xi_0X)$ and that of the left derivative is $\Re\tr(\kappa(x)\Xi_0\kappa(x)^{-1}X)$. Hence the first-order
equations give $\Re\tr(\Xi_0X_i)=0$ and $\Re\tr(\kappa(x)\Xi_0\kappa(x)^{-1}X_i)=0$ for all $i$, and by linearity for all $X$
in the Lie algebra of $N$. The principal symbol of $P^{\mathrm{CV}}_j$ is a non-zero multiple of the coefficient of degree
$2r+k+1-j$ of the characteristic polynomial of $\Xi_0$ (\lref{Formalism.FourierDecayAlt.P.CutMP.SmoothTransferRegularity.operator_capelli_vector}),
so all non-leading coefficients vanish and $\Xi_0$ is nilpotent
(\lref{Formalism.FourierDecayAlt.P.CutMP.SmoothTransferRegularity.nilpotent_of_char_poly}). The principal symbols of
$P^R_j$ and $P^L_j$ are non-zero multiples of the coefficients of the characteristic polynomials of the compressions of $\Xi_0$
to the frame of $z'$, resp.\ of $\kappa(x)\Xi_0\kappa(x)^{-1}$ to the frame of $z$
(\lref{Formalism.FourierDecayAlt.P.CutMP.SmoothTransferRegularity.operator_right_capelli} and \lref{Formalism.FourierDecayAlt.P.CutMP.SmoothTransferRegularity.operator_left_capelli});
their vanishing implies that $\Xi_0$ is orthogonal to the Lie algebra of $U_W(z')$, resp.\ $\kappa(x)\Xi_0\kappa(x)^{-1}$ to
that of $U_W(z)$ (\lref{Formalism.FourierDecayAlt.P.CutMP.SmoothTransferRegularity.compact_zero_holds}). Since
conjugation by $\kappa(x)\in G_{k,r}$ preserves the Lie algebra, non-vanishing, nilpotency and the trace form
(\lref{Formalism.FourierDecayAlt.P.CutMP.SmoothTransferRegularity.conjugation_algebraic_holds}),
$\Xi=\kappa(x)\Xi_0\kappa(x)^{-1}$ has the required properties.
\end{proof}

\subsubsection{The gauge transfer}

\begin{definition}[Transfer statements]\label{pf:sh-curv:gtdef}
Let $P$ be a parameter with $k\ge2$ and $B$ a real place datum which is a place for $P$, with the canonical section. Let
$X$ be an actual-operator datum and $d$ a normal-form datum on the radial space whose level datum $L_d$ is an actual level for
$X$ and a star level, satisfying the normal-form relations, with $X.QQ(Y)=d.QQ(Y)$ for all $Y$ (the \emph{action
hypotheses} of \cref{sec:pf:thmS-carrier:model}). For a frame $e=(e_\rho)$ of the fibres (\cref{pf:thmS-carrier:transferdef}(ii))
let $\mathrm{tr}^R_e$ be the transfer of radial functions to families of functions on $Q_{r,T_T}$ (the radial version of
\cref{pf:thmS-carrier:transferdef}(iii)). We say that $e$ \emph{transfers} an operator $A$ of the radial algebra to an operator
$A^{\mathrm{mod}}$ of the model algebra if $\mathrm{tr}^R_e(A\Phi)=A^{\mathrm{mod}}\mathrm{tr}^R_e(\Phi)$ for all radial
$\Phi$.
\begin{enumerate}[label=(\alph*)]
\item The \emph{general transfer} statement asserts: for all such $(P,B,X,d)$ there is a frame $e$ which transfers
$\zeta^{\pm1}$, $\bar\zeta^{\pm1}$, $G^{\pm1}$, $\nabla$ and $\nabla^*$ of $L_d$, the type multiplications $\kappa(f)$, and the
shifts $S_l$ and $S_l^*$ ($1\le l\le k-1$) of $d$ to the corresponding operators of the model algebra.
\item The \emph{rotation $K$-invariance} statement asserts: for all such $(P,B,X,d)$, every frame $e$ which transfers the type
multiplications and is constant in $\zeta$, every operator $R$ on the radial space commuting with the type multiplications
and acting on radial functions by $(R\Phi)(c,g)=\sum_a(E_{aa}\Phi(c,\cdot))(g)+s\,\Phi(c,g)$ for some $s\in\Z$ (the sum over the
diagonal elements of the last block), and every type $\rho$, there is $m\in\Z$ with
$\mathrm{tr}^R_e(R\Phi)_\rho=\zeta\partial_\zeta\mathrm{tr}^R_e(\Phi)_\rho-\bar\zeta\partial_{\bar\zeta}\mathrm{tr}^R_e(\Phi)_\rho
+m\,\mathrm{tr}^R_e(\Phi)_\rho$ for all $\Phi$.
\end{enumerate}
\formal{Formalism.FourierDecayAlt.P.CutMP.Gauge.GC.GTHyp}
\formal{Formalism.FourierDecayAlt.P.CutMP.Gauge.GC.trR}
\formal{Formalism.FourierDecayAlt.P.CutMP.Gauge.GC.GenTransfer}
\formal{Formalism.FourierDecayAlt.P.CutMP.Gauge.GC.GTRRotK}
\end{definition}

\begin{theorem}[General transfer from rotation invariance]\label{pf:sh-curv:gentransfer}
The rotation $K$-invariance statement implies the general transfer statement.
\formallabel{Formalism.FourierDecayAlt.P.CutMP.Gauge.GC.genTransfer_of_rotK}
\formallabel{Formalism.FourierDecayAlt.P.CutMP.Gauge.GC.genTransfer_of_cut}
\formallabel{Formalism.FourierDecayAlt.P.CutMP.Gauge.GC.gTConn_of_leaves}
\formallabel{Formalism.FourierDecayAlt.P.CutMP.Gauge.GC.gTConn_of_cut}
\formallabel{Formalism.FourierDecayAlt.P.CutMP.Gauge.GC.gTConn_of_rotK}
\end{theorem}

The rotation $K$-invariance statement holds (\lref{Formalism.FourierDecayAlt.P.CutMP.Gauge.GC.rotation_K_invariance_holds},
proved from the matrix form of the rotations, the centre condition, the section, an ordinary differential equation along the
rotation orbits and the Wirtinger derivatives), so the general transfer statement holds; it is used in this form in
\cref{pf:thmS-carrier:openQQ}(T2).

\begin{proof}
The frame is built in three steps, each improving the previous one, and the operators are treated in four groups.
\begin{enumerate}[label=(\roman*)]
\item (Functions) For every frame, the multiplication operators $\zeta^{\pm1},\bar\zeta^{\pm1},G^{\pm1}$ are transferred to the
corresponding multiplication operators of the model
(\lref{Formalism.FourierDecayAlt.P.CutMP.Gauge.GC.GT.gauge_function_holds}): on the section these operators are the
multiplications by $\zeta$, $\bar\zeta$, $|\zeta|^2-T_T$ and their inverses, which do not vanish on the domain.
\item (Types) There is a frame transferring the type multiplications $\kappa(f)$
(\lref{Formalism.FourierDecayAlt.P.CutMP.Gauge.GC.gauge_kappa_holds}, obtained from a pointwise statement and the
$\kappa$-line statement of \cref{sec:pf:thmS-carrier:model}, which provide unit vectors spanning the type components of the
fibres).
\item (Connection) Under the rotation $K$-invariance, every frame as in (ii) can be replaced by a frame which in addition
transfers $\nabla$ and $\nabla^*$ (\lref{Formalism.FourierDecayAlt.P.CutMP.Gauge.GC.gTConn_of_rotK}). This is the
\emph{hermitian curvature} step. By the locality of the gauge
(\lref{Formalism.FourierDecayAlt.P.CutMP.Gauge.GC.gauge_locality_holds}), $\nabla$ and $\nabla^*$ act on the transfer
by a frame $e$ through smooth coefficient functions $A_\rho,B_\rho$ on the annulus $\mathrm{Ann}=\{|\zeta|^2>T_T\}$, one for
each type. One then looks for phases $u_\rho\colon\mathrm{Ann}\to\C$ with $|u_\rho|=1$ and
\[
\bar u_\rho\,\partial_\zeta u_\rho=A^{\mathrm{mod}}_\rho-A_\rho,\qquad
\bar u_\rho\,\partial_{\bar\zeta}u_\rho=-\overline{(A^{\mathrm{mod}}_\rho-A_\rho)},
\]
$A^{\mathrm{mod}}_\rho$ the coefficient of the model. Such phases exist by an integration lemma on the annulus
(\lref{Formalism.FourierDecayAlt.P.CutMP.Gauge.GC.annulus_condition_holds}), applied to $a_\rho=A^{\mathrm{mod}}_\rho-A_\rho$,
whose hypotheses are: $a_\rho$ is smooth, $\Re\partial_{\bar\zeta}a_\rho=0$ on $\mathrm{Ann}$ (the curvature condition,
\lref{Formalism.FourierDecayAlt.P.CutMP.Gauge.GC.gTCCurv_of_ref}), and $\Im\oint_{|w|=R}a_\rho(w)\,dw\in\pi\Z$ for all
$R$ with $R^2>T_T$ (the period condition, \lref{Formalism.FourierDecayAlt.P.CutMP.Gauge.GC.gTCHol_of_ref}); both
conditions are obtained from the curvature computation in a reference frame
(\lref{Formalism.FourierDecayAlt.P.CutMP.Gauge.GC.gTCRef_of_rotK}, where the rotation $K$-invariance is used to control
the trace of the curvature); the reality condition
(\lref{Formalism.FourierDecayAlt.P.CutMP.Gauge.GC.gauge_reality_holds}) relates the coefficients of $\nabla$ and $\nabla^*$,
so that the same phases serve for both (\lref{Formalism.FourierDecayAlt.P.CutMP.Gauge.GC.gTCPhase_of_cut}).
The frame $(u_\rho e_\rho)_\rho$ is again a frame, still transfers the type multiplications, and transfers $\nabla,\nabla^*$
(\lref{Formalism.FourierDecayAlt.P.CutMP.Gauge.GC.gauge_transform_holds}).
\item (Shifts) Every frame as in (iii) can be replaced by a frame which in addition transfers $S_l$ and $S_l^*$ for $1\le
l\le k-1$ (\lref{Formalism.FourierDecayAlt.P.CutMP.Gauge.GC.second_shift_gauge_holds}, proved from the pointwise
adjointness of the second shift, \lref{Formalism.FourierDecayAlt.P.CutMP.Gauge.GC.second_shift_adjoint_pointwise}).
\end{enumerate}
Given $(P,B,X,d)$ with the action hypotheses, (ii) gives a frame $e_0$ transferring the type multiplications, (iii) a frame
$e_1$ transferring them and $\nabla,\nabla^*$, and (iv) a frame $e$ transferring all of these and the shifts; by (i), $e$ also
transfers the multiplication operators. This is the general transfer statement for $(P,B,X,d)$.
\end{proof}

\subsubsection{The head-top identity}

\begin{definition}[Head-top identity and top truncation]\label{pf:sh-curv:htdef}
We use the place packages $p$ of \cref{sec:pf:thmS-carrier:built-defs} (a tower, a base frame, a top lift $\Theta$ and a type
projector) and the notation of \cref{sec:pf:thmS-carrier:leaves}: $RR_p(F)$ are the local kernels of $p$ built from a
coefficient $F$, $\mathrm{slv}(F,w,h)$ the slice vectors of $F$ at a local point $h$, and $\zeta(x)$ the last leading minor of
$x\in\mathrm{Dom}_R$ at the top level.
\begin{enumerate}[label=(\alph*)]
\item The \emph{head-top identity} asserts: for every place package $p$ whose top lift is of class $D^+$ with some constant,
every compactly supported coefficient $F$, cut word $w$, local points $h,h'$ of size $r$, every actual-operator datum at the
top level of the tower, and all radial representatives $\Phi_e,\Phi'_e$ of the slice vectors of $F$ at $h$, $h'$, the top
density of $(\Phi,\Phi')$ for the word $w$ at every $x\in\mathrm{Dom}_R$ equals
$RR_p(F)\bigl(r+1,w,h\cdot(1,\zeta(x)),h'\cdot(1,\zeta(x))\bigr)$.
\item The \emph{top truncation} statement asserts: for every place package $p$ there are smooth compactly supported functions
$b,b'$ on the coordinate space of the top level, a smooth function $\beta$ on the invertible $r\times r$ matrices and
$\kappa\ne0$, with a compactness property for the supports under the upper unipotent group, such that for every function
$\Phi(Z,y)$ which is $U_W(z)\times U_W(z')$-equivariant in $y$ and transforms under upper unitriangular $u$ by the character
$\psi_U(p,u)$ in $Z$, the top pull-back of the truncated coefficient $F_2(b,b',\beta\Phi)$ equals $\kappa\Phi$ at all
$(Z,y)$ with $\det Z\ne0$ and $y$ in the top level subgroup.
\end{enumerate}
\formal{Formalism.FourierDecayAlt.KB.HeadTop}
\formal{Formalism.FourierDecayAlt.P.CutEX.MkTop}
\end{definition}

\begin{theorem}[Head-top identity from top truncation]\label{pf:sh-curv:headtop}
The top truncation statement implies the head-top identity.
\formallabel{Formalism.FourierDecayAlt.P.pb_headTop_of_mkTop}
\formallabel{Formalism.FourierDecayAlt.P.pb_headTop_of}
\formallabel{Formalism.FourierDecayAlt.P.pb_exhPic_of_mkTop}
\end{theorem}

The top truncation statement holds (\lref{Formalism.FourierDecayAlt.P.top_truncation_of_averages}); this gives the
statement (L2) of \cref{pf:thmS-carrier:builtleaves}.

\begin{proof}
The proof passes through an \emph{exhaustion picture} and the positivity of the top unipotent integral.
\begin{enumerate}[label=(\roman*)]
\item (Exhaustion picture) For data as in (a), with an actual-operator datum also at the base of the tower, every cut word
$w$, index $e$ and local point $h$ of size $r$, there are compactly supported coefficients $F_1,F_2$ such that the local vectors
of $\Theta(w,e,F)$ and $\Theta(w,e,F_1)$ at the points $h\cdot(1,\zeta)$ agree, the image of the radial representative of
$F_1$ under the $e$-th component of $\mathbf F_wb_T(w)$ is a radial representative of $F_2$, and the Bessel transform of the
slice of $F_2$ at $h$ is, on the top level, the $e$-th component of $\mathbf F_w$ applied to the radial representatives of the
slice vectors of $F$ (\lref{Formalism.FourierDecayAlt.P.pb_exhPic_of_mkTop}). It is derived from the top truncation
statement and a list of local facts: the exceptional assumption
(\lref{Formalism.FourierDecayAlt.P.CutEX.exceptional_assumption_holds}), the unipotent weight at the top, the exhaustion
of the slices at the top, the smoothness, right action, support and equivalence of pull-backs
(\lref{Formalism.FourierDecayAlt.P.pullback_support}), the compact truncation construction
(\lref{Formalism.FourierDecayAlt.P.compact_truncation_construction}), the existence of top points, a determinant
statement, the exhaustion of head families (\lref{Formalism.FourierDecayAlt.P.head_family_exhaustion}), radial slices
and extensions at exceptional levels, the residue at the top, the regularity of the frame parameters
(\lref{Formalism.FourierDecayAlt.P.frame_parameter_regularity}) and the level-zero picture of slices.
\item (Top unipotent integral) The integral over the highest unipotent subgroup appearing in the top density is positive and
given by an explicit formula (\lref{Formalism.FourierDecayAlt.P.top_unipotent_integral}).
\item (Head weights) The weights of the head of a local point are computed at every point
(\lref{Formalism.FourierDecayAlt.P.head_weight_at}).
\end{enumerate}
From (i)--(iii), the head-top identity follows by radial exhaustion
(\lref{Formalism.FourierDecayAlt.P.head_top_of_radial_exhaustion}): by (i) the top density of the radial representatives
is the Hilbert--Schmidt pairing of the Bessel transforms of the slices of $F_2$ at $h,h'$ on the top level; by (ii) and (iii)
this pairing is the local kernel $RR_p(F)$ at the points $h\cdot(1,\zeta(x))$, $h'\cdot(1,\zeta(x))$.
\end{proof}

\subsubsection{Theorem P for compactly supported test functions}

The last result of this subsection is the input \ref{pf:thmS-carrier:I1} of the carrier statement
(\cref{pf:thmS-carrier:needs}). We state it, construct the subspaces $K_i$ explicitly, and reduce the theorem by a chain of
reductions to analytic properties of the radial operators; the proof is assembled at the end of the subsection. Throughout, $P$, $B$,
$(X,\mu,L,\mathrm{st})$ and $l_0$ are as in \cref{pf:sh-curv:raddef}, with $1\le l_0\le k$, and $S$ is a finite set.

\begin{definition}[Level forms, densities and Theorem P]\label{pf:sh-curv:thmPdef}
Put $I_S=\{0,1\}^S$. A \emph{family} is a map $x\colon I_S\to\mathcal R$; a matrix $M=(M_{ab})_{a,b\in I_S}$ of operators on
$\mathcal R$ acts on families by $(Mx)_a=\sum_bM_{ab}x_b$. The \emph{restriction} of a function $J$ on $\mathrm{Mat}_{k,r}$ is
$\mathrm{res}(J)(c,g)=\prod_{a<r}|c_a|^{m_a-1}J(w(c)g)$, where $w(c)$ is the cell element with coordinates
$c\in(\C^\times)^r$ and the $m_a$ are explicit exponents; $\Phi\in\mathcal R$ is a \emph{radial representative} of $F$ if
$\Phi=\mathrm{res}(J_F)$ on $\mathrm{Dom}_R$.
\begin{enumerate}[label=(\alph*)]
\item (Words) A \emph{letter} is a triple $(\pm,s,Y)$ of a sign, a site $s\in S$ and a point $Y\in\C$. A word $w$ (a list of
letters) is \emph{admissible} if its sites are distinct; its points $Y$ are real, with $Y,Y\pm1\in\mathrm{inv}$ and
$q(Y)q(Y-1)\ne0$; no two points differ by an integer of absolute value at most $2$; every point lies on the cut of $P$ at level
$l_0$, i.e.\ $u_{l_0}+2\le-Y\le u'_{l_0}-2$ or $u'_{l_0}+2\le-Y\le u_{l_0}-2$; and distinct points satisfy $|Y-Y'|>2$.
Through $(L,\mathrm{st})$ a word $w$ defines matrices of operators on $\mathcal R$: the \emph{$RR$-matrices} $RR_i(w)$
($0\le i\le r$), built from the level-$i$ quadratic operators $QQ_i(Y)=t_i(\bar Y)^*\,QQ(Y)\,t_i(Y)$ of its letters ($t_i$ the
product of the transfer matrices of the sites $l\ge i$ of $L$); the \emph{slice matrices} $b_T(l,w)$ ($0\le l<r$), products
of the letter matrices of $w$ at the site $l$; the frame $\mathbf F_w$ of $w$ and its weight $d_w$; and, for $s\in\C$ and
operators $W,W'$ on $\mathcal R$, the \emph{wave matrices} $V_w(s;W,W')$. Put $b_T(w)=b_T(r-1,w)\cdots b_T(0,w)$.
\item (Forms and densities) Let $f_s(t)=\exp(st-\bar s\bar t)$ for $s,t\in\C$, a unitary character of $\C$. For $0\le i\le r$,
$c\in(\C^\times)^r$ and $c'=(c'_l)_{l\ge i}$ let $c\sqcup_ic'$ be $c$ with the coordinates $c_l$, $l\ge i$, replaced by $c'_l$.
For families $x,y$ put
\begin{equation}\label{eq:pf:sh-curv:form}
\mathrm{form}_i(x,y;c)=\int_{\C^{r-i}}\prod_{l\ge i}|c'_l|^{-2}\sum_{a\in I_S}
\int\bigl\langle x_a(c\sqcup_ic',g),y_a(c\sqcup_ic',g)\bigr\rangle_{\mathrm{HS}}\,d\mu(g)\,dc' ,
\end{equation}
and let $RR_i(w;x,y;c)=\mathrm{form}_i(RR_i(w)x,y;c)$ be the \emph{$RR$-form}. For $0\le l<r$ and $t\in\C^\times$ let
$c[l\leftarrow t]$ be $c$ with $c_l$ replaced by $t$; the \emph{phase density} is
$\mathrm{ph}_l(w;x,y;c)(t)=|t|^{-2}RR_{l+1}\bigl(w;b_T(l,w)x,b_T(l,w)y;c[l\leftarrow t]\bigr)$, and the \emph{top density}
is $\mathrm{top}_w(x,y;c,g)=\sum_e\bigl\langle(d_w\mathbf F_wx)_e(c,g),(\mathbf F_wy)_e(c,g)\bigr\rangle_{\mathrm{HS}}$.
\item (Theorem P) Let $K_0,\dots,K_r$ be subspaces of $\mathcal R$; a family is \emph{in $K_i$} if its entries are. $K$
satisfies \emph{Theorem P} if for every admissible word $w$:
\begin{itemize}
\item[(TopP)] for families $x,y$ in $K_r$ and $c\in(\C^\times)^r$, $\mathrm{top}_w(x,y;c,\cdot)$ is $\mu$-integrable,
$\int f_s(\zeta_r(g))\,\mathrm{top}_w(x,y;c,g)\,d\mu(g)=RR_r(w;V_w(s;W,W')x,y;c)$ whenever $W,W'$ are the multiplication
operators by $f_s(\zeta_r)$ and $f_{-s}(\zeta_r)$, and $RR_r(w;x,x;c)\ge0$;
\item[(StepP)] for $0\le l<r$, families $x,y$ in $K_l$ and $c\in(\C^\times)^r$ there is a continuous function $\varphi$ on
$\C$ which equals $\mathrm{ph}_l(w;x,y;c)$ on $\C^\times$, satisfies $|\varphi(t)|\le C(1+|t|)^{-4}$ for some $C$, and
$\int_\C f_s\varphi=RR_l(w;V_w(s;W,W')x,y;c)$ whenever $W,W'$ are the multiplication operators by $f_{\pm s}(c_l)$;
\item[(PosP)] for every $i$ and $c\in(\C^\times)^r$, $(x,y)\mapsto RR_i(w;x,y;c)$ is sesquilinear and non-negative on families
in $K_i$.
\end{itemize}
$K$ satisfies the \emph{slice inclusions} if $b_T(l,w)x$ is in $K_{l+1}$ for all admissible $w$, $l<r$ and families $x$ in
$K_l$, and the \emph{composite level-zero formula} if for all admissible $w$, families $x,y$ in $K_0$ and $c\in(\C^\times)^r$
the function $(c',g)\mapsto\prod_l|c'_l|^{-2}\,\mathrm{top}_w(b_T(w)x,b_T(w)y;c',g)$ is $dc'\,d\mu$-integrable and its
integral is $RR_0(w;x,y;c)$.
\end{enumerate}
\formal{Formalism.FourierDecayAlt.Q.resG}
\formal{Formalism.FourierDecayAlt.P5.WordOK}
\formal{Formalism.FourierDecayAlt.P5.form}
\formal{Formalism.FourierDecayAlt.P5.RRform}
\formal{Formalism.FourierDecayAlt.P5.ph}
\formal{Formalism.FourierDecayAlt.P5.topDens}
\formal{Formalism.FourierDecayAlt.P5.ThmPOn}
\formal{Formalism.FourierDecayAlt.P5.SliceIn}
\formal{Formalism.FourierDecayAlt.P5.CompositeFormulaLevelZero}
\end{definition}

\begin{theorem}[Theorem P for compactly supported test functions]\label{pf:sh-curv:actualcpt}
Let $P$ be a parameter, $B$ a real place datum, $(X,\mu,L,\mathrm{st})$ an actual-operator datum at a level $1\le l_0\le k$, and
$S$ a finite set. There are subspaces $K_0,\dots,K_r$ of $\mathcal R$ which satisfy Theorem P, the slice inclusions and the
composite level-zero formula, and such that every compactly supported $F$ on $\mathrm{Mat}_{k,r}$ which is of class
$\mathrm{CC}(A,m)$ for all $A,m$ (in particular smooth near $G_{k,r}$) and satisfies $G_f\ge\eps$ on $\supp F\cap G_{k,r}$ for
some $\eps>0$ has a radial representative in $K_0$.
\formal{Formalism.FourierDecayAlt.P5.ThmPActualCpt}
\formallabel{Formalism.FourierDecayAlt.P.CutSE.actual_compact_theorem}
\formallabel{Formalism.FourierDecayAlt.P.CutSE.thmPActualCpt_of_R}
\formallabel{Formalism.FourierDecayAlt.P.CutSE.thmPActualCpt_of_spec}
\formallabel{Formalism.FourierDecayAlt.P.CutSE.thmPActualCpt_of_leaves}
\formallabel{Formalism.FourierDecayAlt.P.CutSE.radDecay_of_spec}
\formallabel{Formalism.FourierDecayAlt.P.CutSE.leibniz_class}
\formallabel{Formalism.FourierDecayAlt.P.CutSE.density_smooth_class}
\end{theorem}

\begin{remark}\label{pf:sh-curv:actualcpt-rem}
The formal statement takes in addition a real number $A\ge127(r+1)^3(|S|n^2+n^3)$ and an integer
$m\ge5+3(r+1)(|S|+2)n+|S|r$, $n=2r+k+1$, as hypotheses; they do not occur in the conclusion and are not used in the proof.
For $r=0$ all conditions concerning levels $l<r$ are empty.
\end{remark}

The proof is a chain of reductions: Theorem P is deduced from a list of properties of the top densities and the phase densities
(\cref{pf:sh-curv:thmPsteps}); these are verified for explicit subspaces, the Hermitian spans, except for the decay of the
phase densities (\cref{pf:sh-curv:assembly}); a transport between levels reduces this decay to the level zero
(\cref{pf:sh-curv:levtrans}), where it follows from an expansion of the phase density in derivatives of densities with
rapidly decaying Fourier transforms (\cref{pf:sh-curv:stepexp}); the rapid decay comes from a radial identity
(\cref{pf:sh-curv:raddecay}), which in turn rests on the realisation of the radial operators by differential operators
(\cref{pf:sh-curv:spec}).

\begin{proposition}[Theorem P from step data]\label{pf:sh-curv:thmPsteps}
Let $K_0,\dots,K_r$ be subspaces of $\mathcal R$ and $w$ an admissible word, and assume:
\begin{itemize}
\item[(TI)] for families $x,y$ in $K_r$ and $c\in(\C^\times)^r$, $\mathrm{top}_w(x,y;c,\cdot)$ is $\mu$-integrable and
$RR_r(w;x,y;c)=\int\mathrm{top}_w(x,y;c,g)\,d\mu(g)$;
\item[(TN)] $\mathrm{top}_w(x,x;c,g)\ge0$ for families $x$ in $K_r$ and all $c,g$;
\item[(TT)] if $W,W'$ are the multiplication operators by $f_{\pm s}(\zeta_r)$ and $x,y$ are in $K_r$, then $V_w(s;W,W')x$
is in $K_r$ and $\mathrm{top}_w(V_w(s;W,W')x,y;c,g)=f_s(\zeta_r(g))\,\mathrm{top}_w(x,y;c,g)$;
\item[(SI)] $b_T(l,w)x$ is in $K_{l+1}$ for $l<r$ and families $x$ in $K_l$;
\item[(SD)] for $l<r$, families $x,y$ in $K_l$ and $c\in(\C^\times)^r$ there is a continuous $\varphi$ on $\C$ which equals
$\mathrm{ph}_l(w;x,y;c)$ on $\C^\times$, with $|\varphi(t)|\le C(1+|t|)^{-4}$ for some $C$ and $\int_\C\varphi=RR_l(w;x,y;c)$;
\item[(ST)] if $l<r$, $W,W'$ are the multiplication operators by $f_{\pm s}(c_l)$ and $x,y$ are in $K_l$, then $V_w(s;W,W')x$
is in $K_l$ and $\mathrm{ph}_l(w;V_w(s;W,W')x,y;c)(t)=f_s(t)\,\mathrm{ph}_l(w;x,y;c)(t)$ for $t\ne0$;
\item[(FS)] $(x,y)\mapsto RR_i(w;x,y;c)$ is sesquilinear on families in $K_i$, for all $i$ and $c$.
\end{itemize}
Then (TopP), (StepP) and (PosP) hold for $w$.
\formallabel{Formalism.FourierDecayAlt.P5.thmPFromStepsAll}
\formallabel{Formalism.FourierDecayAlt.P5.nonneg_all_levels}
\end{proposition}

\begin{proof}
\emph{Positivity at all levels.} We show $RR_i(w;x,x;c)\ge0$ for families $x$ in $K_i$ and $c\in(\C^\times)^r$, by descending
induction on $i$. For $i=r$, $RR_r(w;x,x;c)=\int\mathrm{top}_w(x,x;c,g)\,d\mu(g)\ge0$ by (TI) and (TN). Let $i=l<r$ and let
$\varphi$ be as in (SD) for the pair $(x,x)$, so that $RR_l(w;x,x;c)=\int\varphi$. For $t\ne0$ the point $c[l\leftarrow t]$
lies in $(\C^\times)^r$ and $b_T(l,w)x$ is in $K_{l+1}$ by (SI), so
$\varphi(t)=|t|^{-2}RR_{l+1}(w;b_T(l,w)x,b_T(l,w)x;c[l\leftarrow t])\ge0$ by induction. The set of $t$ with
$\varphi(t)\in[0,\infty)$ is closed, since $\varphi$ is continuous, and contains the dense set $\C^\times$; hence $\varphi\ge0$
on $\C$ and $RR_l(w;x,x;c)=\int\varphi\ge0$.

\emph{(TopP).} The integrability is (TI). For $W,W'$ as in (TopP), $V_w(s;W,W')x$ is in $K_r$ by (TT), so by (TI) and (TT)
\begin{align*}
RR_r(w;V_w(s;W,W')x,y;c)&=\int\mathrm{top}_w(V_w(s;W,W')x,y;c,g)\,d\mu(g)\\
&=\int f_s(\zeta_r(g))\,\mathrm{top}_w(x,y;c,g)\,d\mu(g).
\end{align*}
The non-negativity of $RR_r(w;x,x;c)$ is the case $i=r$ above.

\emph{(StepP).} Take $\varphi$ from (SD). Given $s$ and $W,W'$ as in (StepP), let $\psi$ be the function given by (SD) for the
pair $(V_w(s;W,W')x,y)$, which is in $K_l$ by (ST); so $\int\psi=RR_l(w;V_w(s;W,W')x,y;c)$. The continuous functions $f_s\varphi$
and $\psi$ agree on $\C^\times$ by (ST), hence everywhere, and $\int f_s\varphi=\int\psi$.

\emph{(PosP).} Sesquilinearity is (FS); non-negativity was shown above.
\end{proof}

\begin{definition}[Hermitian classes]\label{pf:sh-curv:herdef}
In the situation of \cref{pf:sh-curv:actualcpt} the \emph{Hermitian classes} are the smallest subsets $H_0,\dots,H_r$ of
$\mathcal R$ such that
(i) $H_0$ contains the radial representatives of all $F$ as in \cref{pf:sh-curv:actualcpt};
(ii) $b_T(l,w)_{ab}$ maps $H_l$ into $H_{l+1}$ for all $l<r$, admissible $w$ and $a,b\in I_S$;
(iii) $V_w(s;W,W')_{ab}$ maps $H_l$ into itself for all $l<r$, admissible $w$, $s\in\C$ and $a,b\in I_S$, if $W,W'$ are the
multiplication operators by $f_{\pm s}(c_l)$;
(iv) $V_w(s;W,W')_{ab}$ maps $H_r$ into itself for all admissible $w$, $s$ and $a,b$, if $W,W'$ are the multiplication
operators by $f_{\pm s}(\zeta_r)$.
The \emph{Hermitian spans} are the linear spans $K^{\mathrm{her}}_i$ of the $H_i$.
\formal{Formalism.FourierDecayAlt.P5.HerGenC}
\formal{Formalism.FourierDecayAlt.P5.herKc}
\end{definition}

\begin{proposition}[Theorem P for the Hermitian spans]\label{pf:sh-curv:assembly}
Assume that for all data as in \cref{pf:sh-curv:actualcpt} the Hermitian spans satisfy
\begin{itemize}
\item[(SD)] the step-density property of \cref{pf:sh-curv:thmPsteps} for every admissible word, and
\item[(FI)] for all admissible $w$, $0\le i\le r$, $c\in(\C^\times)^r$ and families $x,y$ in $K^{\mathrm{her}}_i$, the outer
integrand $c'\mapsto\prod_{l\ge i}|c'_l|^{-2}\sum_a\int\langle(RR_i(w)x)_a,y_a\rangle_{\mathrm{HS}}(c\sqcup_ic',g)\,d\mu(g)$ of
$RR_i(w;x,y;c)$ in \eqref{eq:pf:sh-curv:form} is integrable on $\C^{r-i}$.
\end{itemize}
Then \cref{pf:sh-curv:actualcpt} holds, with $K=K^{\mathrm{her}}$.
\formallabel{Formalism.FourierDecayAlt.P5.HerC.thmPActualCpt_of_outer}
\formallabel{Formalism.FourierDecayAlt.P5.HerC.stepFormCptAll_of}
\formallabel{Formalism.FourierDecayAlt.P5.HerC.thmPActualCpt_of_stepForm}
\formallabel{Formalism.FourierDecayAlt.P5.HerC.thmPActualCpt_of_jf}
\formallabel{Formalism.FourierDecayAlt.P5.HerC.thmPActualCpt_of_parts}
\formallabel{Formalism.FourierDecayAlt.P5.HerC.herClassesCpt_of_parts}
\formallabel{Formalism.FourierDecayAlt.P5.HerC.thmPActualCpt_of_herClassesCpt}
\end{proposition}

\begin{proof}
Put $K=K^{\mathrm{her}}$. By construction $K$ satisfies the slice inclusions, is stable under the wave matrices occurring in
(TT) and (ST), and $K_0$ contains the radial representatives of all $F$ as in the theorem
(\lref{Formalism.FourierDecayAlt.P5.HerC.herKc_sliceIn}, \lref{Formalism.FourierDecayAlt.P5.HerC.herKc_waveK}, \lref{Formalism.FourierDecayAlt.P5.HerC.herKc_waveZ}, and \lref{Formalism.FourierDecayAlt.P5.HerC.herKc_jf}).
Every such $F$ has a radial representative: the restriction of $J_F$ to $\mathrm{Dom}_R$, extended by zero, lies in
$\mathcal R$ (\lref{Formalism.FourierDecayAlt.P5.radJFExists}). This is the last assertion of the theorem. Fix an
admissible word $w$; we verify the hypotheses of \cref{pf:sh-curv:thmPsteps} and the composite level-zero formula.

(TI) Every $\Phi\in K_r$ has compact support in $g$: for every $c\in(\C^\times)^r$, $\Phi(c,\cdot)$ vanishes outside a compact
subset of $G_{k,r}$. Indeed, radial representatives of compactly supported $F$ have this property locally uniformly in $c$
(\lref{Formalism.FourierDecayAlt.P5.HerC.jfCptLocAll}, which uses that the leading minors of $\widehat Z$ are invariant
under the action of $\mathcal S$ through which $J_F$ is defined), the functions with this local property form a subspace
stable under the slice and wave matrices (\lref{Formalism.FourierDecayAlt.P5.HerC.herKc_le_lc}), and the local property
implies compact support at every $c$ (\lref{Formalism.FourierDecayAlt.P5.HerC.herCptTopAll_of_jf}). For such families the
operators of the frame can be moved across the inner product,
$\mathrm{form}_r(\mathbf F_w^*d_w\mathbf F_wx,y;c)=\int\mathrm{top}_w(x,y;c,g)\,d\mu(g)$ with an integrable density
(\lref{Formalism.FourierDecayAlt.P5.HerA.adjTop_cpt}), and $\mathbf F_w^*d_w\mathbf F_w=RR_r(w)$ for admissible words
(\lref{Formalism.FourierDecayAlt.P.frame_of_admissible_word} and \lref{Formalism.FourierDecayAlt.P5.ofWordRR}); this is (TI)
(\lref{Formalism.FourierDecayAlt.P5.topIdentFromParts}).

(TN) By the normal form of the actual operators (\lref{Formalism.FourierDecayAlt.P.normal_form_actual}) there are type
multiplications $\kappa$ with $x(Y)=\kappa\bigl((x_0(\rho,Y))_\rho\bigr)$ and $\xi(Y)=\kappa\bigl((x_0(\rho,Y)^{-1})_\rho\bigr)$
for $Y\in\mathrm{inv}$, and $\kappa(f)$ is positive on the fibres when $f$ is positive
(\lref{Formalism.FourierDecayAlt.P5.xKappaOfNF} and \lref{Formalism.FourierDecayAlt.P5.kappaFibrePos}). Hence the weight of an
admissible word is $d_w=(\zeta\bar\zeta G^{-1})^N\,\mathrm{diag}_e\,\kappa(c_e)$ with explicit functions $c_e>0$ (products of
values of $x_0(\rho,\cdot)$ and $q$ at the points of $w$, positive on the cut by interlacing), and
$|\zeta|^2G^{-1}\ge0$ on $\mathrm{Dom}_R$; so every term of $\mathrm{top}_w(x,x;c,g)$ is non-negative
(\lref{Formalism.FourierDecayAlt.P5.topDensNonnegOfKappa}).

(TT) The wave matrix at the top multiplies the top density by $f_s(\zeta_r)$
(\lref{Formalism.FourierDecayAlt.P5.topTwistDensOfZeta}, using the smoothness of $\zeta_r$ on the radial domain,
\lref{Formalism.FourierDecayAlt.P.actual_zeta_smooth_alt}), and $K_r$ is stable under it by construction.

(SI) holds by construction, and (SD) is the hypothesis.

(ST) follows from the plane-wave vector identity (\lref{Formalism.FourierDecayAlt.P.plane_wave_vector_identity}), the
commutation of the multiplication operators by $f_{\pm s}(c_l)$ with the site operators and with the quadratic operators
(\lref{Formalism.FourierDecayAlt.P5.waveSiteAll} and \lref{Formalism.FourierDecayAlt.P5.waveCommQQAll}) and the stability of
$K_l$ (\lref{Formalism.FourierDecayAlt.P5.stepTwistFromParts}).

(FS) The inner integrals in \eqref{eq:pf:sh-curv:form} exist for families in $K_i$
(\lref{Formalism.FourierDecayAlt.P5.HerC.formIntG_herKc}) and the outer integrand is integrable by (FI); so
$RR_i(w;x,y;c)$ is an integral of an expression sesquilinear in $(x,y)$, and linearity of the integral gives (FS)
(\lref{Formalism.FourierDecayAlt.P5.formSesqOfInt}).

By \cref{pf:sh-curv:thmPsteps}, $K$ satisfies Theorem P. Finally, the composite level-zero formula is proved by descending
induction on $i$: with $b_{\ge i}=b_T(r-1,w)\cdots b_T(i,w)$,
\[
RR_i(w;x,y;c)=\int\prod_{l\ge i}|c'_l|^{-2}\,\mathrm{top}_w(b_{\ge i}x,b_{\ge i}y;c\sqcup_ic',g)\,dc'\,d\mu(g),
\]
with integrable integrand. The case $i=r$ is (TI), and the step from $i+1$ to $i$ splits off the coordinate $c'_i$ (a
measure-preserving rearrangement of the product space) and uses (SD), (SI), the integrability of the step functions
$\varphi$, which follows from their quartic decay, and the measurability and non-negativity of the top densities
(\lref{Formalism.FourierDecayAlt.P5.integral_formula_of_steps} and \lref{Formalism.FourierDecayAlt.P5.topDensMeasAll}). For
$i=0$ this is the composite level-zero formula.
\end{proof}

\begin{proposition}[Reduction to level zero]\label{pf:sh-curv:levtrans}
Say that a function $D$ on $\C$ has \emph{decay data of order $n$} if $D$ and $(1+|t|^2)^2D$ are integrable, $D$ is continuous
on $\C^\times$, and $\int f_sD$ and $\int f_s\,(1+|t|^2)^2D$ are $O\bigl((1+|s|)^{-(n+3)}\bigr)$ uniformly in $s\in\C$. Assume
\begin{itemize}
\item[(E0)] for all data as in \cref{pf:sh-curv:actualcpt} with $r\ge1$, admissible $w$, families $x,y$ in
$K^{\mathrm{her}}_0$ and $c\in(\C^\times)^r$ there are $n\in\N$, functions $D_1,\dots,D_K$ with decay data of order $n$,
integers $j_i\le n$, vectors $v_i\in\C^{j_i}$ and numbers $a_i$ such that
\[
\mathrm{ph}_0(w;x,y;c)(t)=\sum_ia_i\,D^{j_i}D_i(t)[v_i]\quad(t\ne0),\qquad\sum_{i:\,j_i=0}a_i\int_\C D_i=RR_0(w;x,y;c),
\]
where $D^jD(t)[v]$ is the $j$-th real derivative at $t$ evaluated at $v$.
\end{itemize}
Then the hypotheses (SD) and (FI) of \cref{pf:sh-curv:assembly} hold; hence \cref{pf:sh-curv:actualcpt} holds.
\formallabel{Formalism.FourierDecayAlt.P5.HerD.thmPActualCpt_of_stepExp0}
\formallabel{Formalism.FourierDecayAlt.P5.HerD.theorem_P_compact_of_step_density}
\formallabel{Formalism.FourierDecayAlt.P5.HerD.stepDensCptAll_of_levT}
\formallabel{Formalism.FourierDecayAlt.P5.HerD.stepDens_of_exp}
\formallabel{Formalism.FourierDecayAlt.P5.HerD.ext_of_derivCombInt}
\formallabel{Formalism.FourierDecayAlt.P5.HerD.derivCombInt_congr}
\end{proposition}

\begin{proof}
\emph{Step 1: from an expansion to a step function.} Let $f$ be a function on $\C^\times$ with an expansion
$f(t)=\sum_ia_iD^{j_i}D_i(t)[v_i]$ as in (E0), with value $I=\sum_{j_i=0}a_i\int D_i$. If $D$ has decay data of order $n$, there
is a $C^n$ function $D'$ on $\C$ which equals $D$ on $\C^\times$, with $\int D'=\int D$, whose derivatives of order $\le n$ are
bounded by $C'(1+|t|)^{-4}$, and whose derivatives of order $1,\dots,n$ are integrable with integral $0$
(\lref{Formalism.FourierDecayAlt.P5.HerC.Fou.smooth_extension_of_fourier_decay}, by Fourier inversion: the transforms of
$D$ and of $(1+|t|^2)^2D$ are $O((1+|s|)^{-(n+3)})$, and $|s|^j(1+|s|)^{-(n+3)}$ is integrable on $\C=\R^2$ for $j\le n$).
Put $\varphi=\sum_ia_iD^{j_i}D'_i[v_i]$. It is continuous; it equals $f$ on $\C^\times$, because $D'_i=D_i$ on the open set
$\C^\times$; $|\varphi(t)|\le\sum_i|a_i|C'_i\prod_m|v_{i,m}|\,(1+|t|)^{-4}$; and $\int\varphi=\sum_{j_i=0}a_i\int D_i=I$, since
the terms with $j_i\ge1$ integrate to $0$. Multiplying the $a_i$ by a constant $\lambda$ gives an expansion of $\lambda f$ with
value $\lambda I$.

\emph{Step 2: level transport.} By \lref{Formalism.FourierDecayAlt.P.level_transport_hermitian} (obtained from the good
behaviour of the slice families, \lref{Formalism.FourierDecayAlt.P.slice_family_good}): for every level $i\le r$ and every
$c$ with $c_l\ne0$ for $l<i$ there are data as in the theorem for the parameter $(k,r-i,T,T',u,u')$ and a linear map
$\mathcal T$ from $\mathcal R$ to the new radial space, $(\mathcal T\Phi)(c'',g)=\Phi(c\sqcup_ic'',\iota(g))$ with the standard
embedding $\iota$ of the smaller unitary group, such that
(a) $\int\langle\mathcal T\Phi,\mathcal T\Psi\rangle_{\mathrm{HS}}(c'',g)\,d\mu_i(g)=\kappa\int\langle\Phi,\Psi\rangle_{\mathrm{HS}}
(c\sqcup_ic'',g)\,d\mu(g)$ for some $\kappa>0$;
(b) $\mathcal T$ intertwines the quadratic, slice and $RR$-matrices at the levels $n+i$ with those at the levels $n$ of the new
data;
(c) admissible words remain admissible;
(d) $\mathcal T$ maps $K^{\mathrm{her}}_i$ into the level-zero Hermitian span of the new data.
By (a) and (b), $RR_{n+i}(w;x,y;c')=\kappa^{-1}RR^{\mathrm{new}}_n(w;\mathcal Tx,\mathcal Ty;c'_{\ge i})$ for all $c'$ which agree
with $c$ below $i$.

\emph{Step 3: (SD).} Let $l<r$, $x,y$ be families in $K^{\mathrm{her}}_l$ and $c\in(\C^\times)^r$, and apply Step 2 with
$i=l$; the new data have $r-l\ge1$. By Step 2, $RR_l(w;x,y;c)=\kappa^{-1}RR^{\mathrm{new}}_0(w;\mathcal Tx,\mathcal Ty;c_{\ge l})$
and $\mathrm{ph}_l(w;x,y;c)(t)=\kappa^{-1}\mathrm{ph}^{\mathrm{new}}_0(w;\mathcal Tx,\mathcal Ty;c_{\ge l})(t)$ for $t\ne0$; the
families $\mathcal Tx,\mathcal Ty$ lie in the level-zero Hermitian span of the new data, and $w$ is admissible for them. By (E0)
for the new data and Step 1 (with $\lambda=\kappa^{-1}$), $\mathrm{ph}_l(w;x,y;c)$ has an expansion with value
$RR_l(w;x,y;c)$, and Step 1 gives the function $\varphi$ required in (SD).

\emph{Step 4: (FI)} is deduced from the level transport as well
(\lref{Formalism.FourierDecayAlt.P5.HerD.formIntOuterCpt_of_levT}). The theorem then follows from
\cref{pf:sh-curv:assembly}.
\end{proof}

\begin{definition}[Analytic inputs at level zero]\label{pf:sh-curv:anadef}
Let the data be as in \cref{pf:sh-curv:actualcpt} with $r\ge1$. The \emph{level-zero class} $E\subseteq\mathcal R$ is the span
of the images of the elements of $K^{\mathrm{her}}_0$ and of $RR_0(w)_{ab}K^{\mathrm{her}}_0$ ($w$ admissible) under all
products of the actual operators of the site $0$: the radial derivatives $p_0,\bar p_0$ and the multiplications by
$c_0^{\pm1},\bar c_0^{\pm1}$. For families $u,v$ and $c\in(\C^\times)^r$ the \emph{density} is
$\mathrm{Dens}(u,v;c)(t)=|t|^{-2}\mathrm{form}_1(u,v;c[0\leftarrow t])$, $t\in\C^\times$; $\partial_t,\partial_{\bar t}$ are the
Wirtinger derivatives. A function $G$ on $\mathrm{Mat}_{k,r}$ is \emph{nice} if it is smooth near $G_{k,r}$, of class
$\mathrm{CC}(A,m)$ for all $A,m$, and satisfies $G_f\ge\eps$ on its support in $G_{k,r}$ for some $\eps>0$. The analytic
inputs are, for families $u,v$ in $E$ and $c\in(\C^\times)^r$:
\begin{itemize}
\item[(LB)] (Leibniz rule) for $t\ne0$,
$\partial_t\mathrm{form}_1(u,v;c[0\leftarrow t])=\mathrm{form}_1(c_0^{-1}p_0u,v;c[0\leftarrow t])+\mathrm{form}_1(u,\bar c_0^{-1}
\bar p_0v;c[0\leftarrow t])$, and the same with $\partial_{\bar t}$ and the roles of $c_0^{-1}p_0$ and $\bar c_0^{-1}\bar p_0$
exchanged;
\item[(DS)] (smoothness) $\mathrm{Dens}(u,v;c)$ is $C^\infty$ on $\C^\times$, and the entries of $u,v$ are square integrable for
$\mathrm{form}_1$ at the points $c[0\leftarrow t]$, $t\ne0$;
\item[(OL)] (one-letter density formula) for every letter $\ell$ with real point there are $e\in\C$,
$\partial\in\{\partial_t,\partial_{\bar t}\}$ and $i,j\in\{0,1\}$ such that, for $t\ne0$,
$\mathrm{Dens}(\check v_0(\ell)u,v_0(\ell)v;c)(t)=\mathrm{Dens}(u,v;c)(t)+e\,\partial\bigl(\mathrm{Dens}(u,E^{ij}_{\ell}v;c)\bigr)(t)$,
where $v_0(\ell),\check v_0(\ell)$ are the two letter matrices of $\ell$ at the site $0$ and $E^{ij}_\ell$ is the elementary
matrix at the site of $\ell$;
\item[(RD)] (radial decay) for nice $G,H$ with radial representatives $\Phi,\Psi$ and every $n\in\N$ there is $C$ with
$|\mathrm{form}_0(W_s\Phi,\Psi;c)|\le C(1+|s|)^{-(n+3)}$ for all $s\in\C$, $W_s$ the multiplication operator by $f_s(c_0)$;
here $\mathrm{form}_0(\Phi,\Psi;c)$, for single elements of $\mathcal R$, is defined by \eqref{eq:pf:sh-curv:form} without the
sum over $a$.
\end{itemize}
\formal{Formalism.FourierDecayAlt.P.CutSE.ClsE}
\formal{Formalism.FourierDecayAlt.P5.HerD.Dens}
\formal{Formalism.FourierDecayAlt.P.CutSE.LeibnizCls}
\formal{Formalism.FourierDecayAlt.P.CutSE.DensSmoothCls}
\formal{Formalism.FourierDecayAlt.P.CutSE.OneLetterDens}
\formal{Formalism.FourierDecayAlt.P.CutSE.RadDecay}
\end{definition}

\begin{proposition}[Derivative expansion at level zero]\label{pf:sh-curv:stepexp}
If (OL), (DS) and (RD) hold for all data with $r\ge1$, then (E0) of \cref{pf:sh-curv:levtrans} holds.
\formallabel{Formalism.FourierDecayAlt.P.CutSE.stepExp0_of_leaves}
\end{proposition}

\begin{proof}
This is \lref{Formalism.FourierDecayAlt.P.CutSE.stepExp0_of_open}, applied with the decay statements described in (ii).
(i) \emph{Letter expansion.} Note $\mathrm{ph}_0(w;x,y;c)=\mathrm{Dens}(RR_1(w)b_T(0,w)x,b_T(0,w)y;c)$. Removing the letters of
$w$ one at a time with (OL), the class $E$ being stable under the letter and elementary matrices, one writes on $\C^\times$
\[
\mathrm{ph}_0(w;x,y;c)=\sum_ia_i\,\partial^{\beta_i}\mathrm{Dens}(RR_0(w)x,v_i;c)
\]
with iterated Wirtinger derivatives $\partial^{\beta_i}$ of order $\le|w|$, families $v_i$ in $E$, and
\[
\sum_{i:\,\beta_i=\emptyset}a_i\,\mathrm{form}_0(RR_0(w)x,v_i;c)=RR_0(w;x,y;c).
\]
Iterated Wirtinger derivatives are linear
combinations of iterated real derivatives, and $\int\mathrm{Dens}(u,v;c)=\mathrm{form}_0(u,v;c)$ by Fubini's theorem; this gives
the expansion required in (E0) (\lref{Formalism.FourierDecayAlt.P.CutSE.expansion_l2_of_open_estimates}).
(ii) \emph{Decay data.} Each density $\mathrm{Dens}(RR_0(w)x,v;c)$ has decay data of every order: it is continuous on
$\C^\times$ by (DS) (\lref{Formalism.FourierDecayAlt.P.CutSE.decayPair2_of_smooth}); its Fourier transform is
$\int f_s(t)\mathrm{Dens}(u,v;c)(t)\,dt=\mathrm{form}_0(W_su,v;c)$ (the outgoing-wave identity,
\lref{Formalism.FourierDecayAlt.P.CutSE.wave_outgoing_orbital}), which decays faster than any power of $1+|s|$ by (RD)
extended to the class (\lref{Formalism.FourierDecayAlt.P.CutSE.clsDecayJ_of}); the weight $(1+|t|^2)^2$ is produced by the
site multiplication operators (\lref{Formalism.FourierDecayAlt.P.CutSE.site_multiplication_orbital}), which gives the same
decay for $(1+|t|^2)^2\mathrm{Dens}$ (\lref{Formalism.FourierDecayAlt.P.CutSE.decayPair3_of} and \lref{Formalism.FourierDecayAlt.P.CutSE.decayPair4_of});
the integrability of the densities, with and without the weight, is established together with the expansion (i).
\end{proof}

\begin{definition}[Wave twists and the radial identity]\label{pf:sh-curv:radiddef}
Let $B$ be a real place datum with $r\ge1$. For $s\in\C$ and a function $F$ on $\mathrm{Mat}_{k,r}$ the \emph{wave twist} is
$w_sF(g)=f_s(\Delta_1(g))F(g)$, $\Delta_1(g)$ the first leading minor of $\widehat Z(g)$. For functions $F,F'$ let
$\beta(F,F')$ be the pairing of the Bessel transforms $J_F,J_{F'}$ over the section of $B$. A \emph{good operator} is a
differential operator $\Theta$ on $\mathrm{Mat}_{k,r}$ with smooth coefficients on the positivity domain of $B$, commuting with
the action of $\mathcal S$, whose coefficients have polynomial growth on $\{G_f\ge\eps\}$ for every $\eps>0$, and such that
$\Theta$ and its formal adjoint $\Theta^*$ map nice functions to nice functions. The \emph{radial identity} holds for good
operators $\Theta_0,\dots,\Theta_3$ if for all $s\in\C$, nice $F$ and $g$ in the image of the section of $B$
\begin{equation}\label{eq:pf:sh-curv:radid}
s\,J_{w_sF}(g)=J_{\Theta_0^*(w_s\Theta_1F)}(g)-J_{\Theta_2^*(w_s\Theta_3F)}(g).
\end{equation}
\formal{Formalism.FourierDecayAlt.P.CutSE.wvG}
\formal{Formalism.FourierDecayAlt.B.beta}
\formal{Formalism.FourierDecayAlt.P.CutSE.RDOpsProp}
\formal{Formalism.FourierDecayAlt.P.CutSE.RDRadIdPropQ}
\end{definition}

\begin{proposition}[Rapid decay from the radial identity]\label{pf:sh-curv:raddecay}
(a) If the radial identity holds for good operators $\Theta_0,\dots,\Theta_3$, then for nice $F,F'$ and every $n\in\N$ there
is $C$ with $|\beta(w_sF,F')|\le C(1+|s|)^{-(n+3)}$ for all $s\in\C$.
(b) If for all data with $r\ge1$ the radial identity holds for some good operators, then (RD) holds.
\formallabel{Formalism.FourierDecayAlt.P.CutSE.radial_identity_of_quadratic}
\formallabel{Formalism.FourierDecayAlt.P.CutSE.rd_iter_pow}
\formallabel{Formalism.FourierDecayAlt.P.CutSE.rd_pow_le}
\formallabel{Formalism.FourierDecayAlt.P.CutSE.rd_decay}
\formallabel{Formalism.FourierDecayAlt.P.CutSE.radDecay_of_cut}
\formallabel{Formalism.FourierDecayAlt.P.CutSE.radDecay_of_ABQ}
\end{proposition}

\begin{proof}
(a) \emph{The identity for $\beta$.} Since $\beta(F,F')$ is computed from the values of $J_F$ and $J_{F'}$ on the image of the
section and is linear in $J_F$, \eqref{eq:pf:sh-curv:radid} gives
$s\,\beta(w_sF,F')=\beta(\Theta_0^*(w_s\Theta_1F),F')-\beta(\Theta_2^*(w_s\Theta_3F),F')$. Good operators can be moved across
$\beta$: $\beta(\Theta^*G,F')=\beta(G,\Theta F')$ for nice $G,F'$ (\lref{Formalism.FourierDecayAlt.P.b_any_of_xi}, which
uses the smoothness, the equivariance and the polynomial growth of the coefficients and the support condition
$G_f\ge\eps$), and wave twists of nice functions are nice
(\lref{Formalism.FourierDecayAlt.P.CutSE.radial_wave_identity}). Hence
\begin{equation}\label{eq:pf:sh-curv:betaid}
s\,\beta(w_sF,F')=\beta(w_s\Theta_1F,\Theta_0F')-\beta(w_s\Theta_3F,\Theta_2F')\qquad(s\in\C,\ F,F'\text{ nice}).
\end{equation}
\emph{Iteration.} We show by induction on $M$ that for nice $F,F'$ there is $K_M(F,F')$ with
$|s|^M|\beta(w_sF,F')|\le K_M(F,F')$ for all $s$. For $M=0$ this is the boundedness of $s\mapsto\beta(w_sF,F')$
(\lref{Formalism.FourierDecayAlt.P.CutSE.radial_boundedness}; recall $|f_s|=1$). If it holds for $M$, then by
\eqref{eq:pf:sh-curv:betaid}
\begin{align*}
|s|^{M+1}|\beta(w_sF,F')|&\le|s|^M|\beta(w_s\Theta_1F,\Theta_0F')|+|s|^M|\beta(w_s\Theta_3F,\Theta_2F')|\\
&\le K_M(\Theta_1F,\Theta_0F')+K_M(\Theta_3F,\Theta_2F'),
\end{align*}
the functions $\Theta_jF,\Theta_jF'$ being nice. \emph{Conclusion.} For $x\ge0$ and $N\in\N$, $(1+x)^N\le2^N(1+x^N)$ (if $x\le1$
the left side is $\le2^N$; otherwise it is $\le(2x)^N$). With $N=n+3$,
$(1+|s|)^{n+3}|\beta(w_sF,F')|\le2^{n+3}\bigl(K_0(F,F')+K_{n+3}(F,F')\bigr)$, which is (a).

(b) Let $G,H$ be nice with radial representatives $\Phi,\Psi$. On the radial domain the first leading minor is the coordinate
$c_0$, so $W_s\Phi$ is a radial representative of the nice function $w_sG$
(\lref{Formalism.FourierDecayAlt.P.CutSE.radial_wave_identity}); and there is a constant $C_1\ne0$, depending only on $B$
and $\mu$, with $\mathrm{form}_0(\Phi',\Psi';c)=C_1\beta(G',H')$ for all $c$ and all radial representatives $\Phi',\Psi'$ of
functions $G',H'$ satisfying all growth bounds (\lref{Formalism.FourierDecayAlt.P.CutSE.radial_beta_identity}). Hence
$|\mathrm{form}_0(W_s\Phi,\Psi;c)|=|C_1|\,|\beta(w_sG,H)|\le|C_1|C(1+|s|)^{-(n+3)}$ by (a).
\end{proof}

\begin{proposition}[Spectral realisation of the radial operators]\label{pf:sh-curv:spec}
Let the data be as in \cref{pf:sh-curv:actualcpt} with $r\ge1$. Write $R(Y)$, $R^\vee(Y)$ for the radial operators of the
level datum built from the entries of $QQ(Y)$, $G_R^{-1}$ for the product of the inverse radial operators and $G^{-1}$, and
$\widehat A(Y)$, $\widehat A^\vee(Y)$ for the operators $\widehat A$ of the level datum. An operator $a$ on $\mathcal R$ is
\emph{realised} by a good operator $\Theta$ if $\mathrm{res}(\Theta J)=a\Phi$ on $\mathrm{Dom}_R$ for every $J$ in the Bessel
class of $(z,z')$ and every $\Phi\in\mathcal R$ with $\Phi=\mathrm{res}(J)$ on $\mathrm{Dom}_R$.
\begin{enumerate}[label=(\alph*)]
\item For every real $Y$ with $Y-1,Y,Y+1\in\mathrm{inv}$, the operators $G_R^{-1}\xi(Y)R(Y)$ and
$G_R^{-1}\xi(Y+1)R^\vee(Y+1)$ are realised by good operators.
\item There are a real $Y$ with $Y,Y+1\in\mathrm{inv}$ and good operators $\Theta_0,\dots,\Theta_3$ realising
$\widehat A^\vee(Y+1)$, $G_R^{-1}\xi(Y)R(Y)$, $G_R^{-1}\xi(Y+1)R^\vee(Y+1)$, $\widehat A(Y)$, whose formal adjoints realise the
adjoint components of these operators.
\item For $Y$ and $\Theta_0,\dots,\Theta_3$ as in (b), the radial identity holds.
\end{enumerate}
\formallabel{Formalism.FourierDecayAlt.P.CutSE.spectral_operator_R}
\formallabel{Formalism.FourierDecayAlt.P.CutSE.spectral_operator_R_adjoint}
\formallabel{Formalism.FourierDecayAlt.P.CutSE.rDSOpR_of_parts}
\formallabel{Formalism.FourierDecayAlt.P.CutSE.rDSOpRv_of_parts}
\formallabel{Formalism.FourierDecayAlt.P.CutSE.rDIdentSpec_of_R}
\formallabel{Formalism.FourierDecayAlt.P.CutSE.rDIdentABQ_of_spec}
\end{proposition}

\begin{proof}
(a) Let $Y$ be real with $Y-1,Y,Y+1\in\mathrm{inv}$. (i) There is $\sigma\in\C$ such that $R(Y)$ acts on radial representatives
as $\sigma$ times an explicit group operator $\mathrm{rh}(Y)$ acting on restrictions
(\lref{Formalism.FourierDecayAlt.P.CutSE.spectral_right_rho_action}); (ii) $\sigma\,\mathrm{rh}(Y)$ is induced on
restrictions by a good operator $\Theta'$ (\lref{Formalism.FourierDecayAlt.P.CutSE.spectral_right_group_action});
(iii) $G_R^{-1}\xi(Y)$ acts on $\mathcal R$ as multiplication by a smooth coefficient $T_0$ of polynomial growth
(\lref{Formalism.FourierDecayAlt.P.CutSE.spectral_operator_X}); (iv) the composite $T_0\Theta'$ is again good
(\lref{Formalism.FourierDecayAlt.P.CutSE.spectral_composition_zero}). For $J$ in the Bessel class and $\Phi=\mathrm{res}(J)$,
\[
G_R^{-1}\xi(Y)R(Y)\Phi=T_0\cdot\bigl(\sigma\,\mathrm{rh}(Y)\mathrm{res}(J)\bigr)=T_0\cdot\mathrm{res}(\Theta'J)=\mathrm{res}(T_0\Theta'J)
\]
on $\mathrm{Dom}_R$, so $T_0\Theta'$ realises $G_R^{-1}\xi(Y)R(Y)$. The same argument with
\lref{Formalism.FourierDecayAlt.P.CutSE.spectral_adjoint_rho_action} and \lref{Formalism.FourierDecayAlt.P.CutSE.spectral_adjoint_right_group}
in (i), (ii), at the point $Y+1$, realises $G_R^{-1}\xi(Y+1)R^\vee(Y+1)$.

(b) This is \lref{Formalism.FourierDecayAlt.P.CutSE.rDIdentSpec_of_open}, applied to (a) and to the following spectral
facts: the realisations of $\widehat A(Y)$ and $\widehat A^\vee(Y)$
(\lref{Formalism.FourierDecayAlt.P.CutSE.spectral_operator_A} and \lref{Formalism.FourierDecayAlt.P.CutSE.spectral_operator_A_adjoint}),
the adjointness of the coordinate pairings and of the quadratic operators
(\lref{Formalism.FourierDecayAlt.P.CutSE.spectral_pairing_adjoint} and \lref{Formalism.FourierDecayAlt.P.CutSE.spectral_quadratic_adjoint}),
the stability of the radial test functions under $\xi(Y)$ (\lref{Formalism.FourierDecayAlt.P.CutSE.spectral_xi_test}) and
the pairing identity (\lref{Formalism.FourierDecayAlt.P.CutSE.spectral_pairing_identity}). The point $Y$ is chosen real and
outside the finite set of shifts of zeros of the type polynomials, so that the relevant shifts of $Y$ lie in $\mathrm{inv}$
(\lref{Formalism.FourierDecayAlt.P.CutSE.spectral_Y_real}).

(c) For $Y$ and $\Theta$ as in (b), the radial identity \eqref{eq:pf:sh-curv:radid} holds
(\lref{Formalism.FourierDecayAlt.P.CutSE.radial_identity_of_spectral_data}): an algebraic identity between the operators
$\theta(Y)=\bigl(\widehat A^\vee(Y+1),G_R^{-1}\xi(Y)R(Y),G_R^{-1}\xi(Y+1)R^\vee(Y+1),\widehat A(Y)\bigr)$ of the level datum
(obtained from the collapse, gauge and rearrangement identities of the radial operators) is transported to the group through
the realisations of (b), the wave identities and the residue formulas of the radial operators.
\end{proof}

\begin{proof}[Proof of \cref{pf:sh-curv:actualcpt}]
If $r=0$, (E0) is empty, and the theorem follows from \cref{pf:sh-curv:levtrans}. Let $r\ge1$. By
\cref{pf:sh-curv:spec} the radial identity holds for good operators, so (RD) holds by
\cref{pf:sh-curv:raddecay}(b). The Leibniz rule (LB) and the smoothness (DS) follow from the radial locality of radial
representatives (\lref{Formalism.FourierDecayAlt.P.CutSE.spectral_radial_locality}: for radial representatives of nice
compactly supported functions, all iterated radial derivatives by $p_0,\bar p_0$ are square integrable for $\mathrm{form}_1$,
with norms locally bounded in $c_0\ne0$) by differentiation under the integral sign
(\lref{Formalism.FourierDecayAlt.P.CutSE.leibnizCls_of_radLoc} and \lref{Formalism.FourierDecayAlt.P.CutSE.densSmoothCls_of_radLoc});
the one-letter density formula (OL) is a computation of the action of a single letter on the density from (LB) and (DS)
(\lref{Formalism.FourierDecayAlt.P.CutSE.one_letter_density_estimate}). By \cref{pf:sh-curv:stepexp}, (E0) holds, and
\cref{pf:sh-curv:levtrans} gives the theorem.
\end{proof}
